\documentclass[11pt]{article}

\usepackage[margin=1in]{geometry}
\usepackage{amsmath}
\usepackage{amssymb}
\usepackage{amsthm}
\usepackage{mathtools}
\usepackage{graphicx}
\usepackage{booktabs}
\usepackage{array}
\usepackage{longtable}
\usepackage[utf8]{inputenc}
\usepackage[T1]{fontenc}
\usepackage{lmodern}
\usepackage{microtype}
\usepackage{xcolor}
\usepackage{listings}
\usepackage[hidelinks]{hyperref}

\lstset{
  basicstyle=\ttfamily\footnotesize,
  breaklines=true,
  showstringspaces=false,
  columns=fullflexible,
}

% Stretch interword spaces in last-resort line breaking so long compound
% \texttt{} identifiers (file paths, function names) don't blow the right
% margin. Standard "long-URL" escape hatch.
\emergencystretch=3em
\setlength{\hfuzz}{0.5pt}
% Allow line breaks at slashes and underscores inside \texttt{}.
\makeatletter
\g@addto@macro\UrlBreaks{\do\/\do\_\do\-}
\makeatother
\newcommand{\stub}[1]{\textbf{[STUB: #1]}}
\newcommand{\Sm}{S(\mathbf{q})}
\newcommand{\bofq}{b(\mathbf{q})}

\title{Liquidity Party\\
       \large A Multi-Asset Logarithmic Market Scoring Rule AMM\\
       \large Deployment Whitepaper}
\author{}
\date{June 11, 2026}

\begin{document}
\maketitle

\begin{abstract}
Liquidity Party is a permissioned-deploy, non-upgradeable multi-asset
automated market maker. The pricing kernel is the Logarithmic Market
Scoring Rule (LMSR) with the liquidity parameter quasi-statically tied
to pool size as $\bofq = \kappa\,\Sm$, but held \emph{fixed for the
duration of each block} and anchored to $\min(\sigma_{\mathrm{swap}},
\sigma_{\mathrm{live}})$, where $\sigma_{\mathrm{swap}}$ is a
block-stepped EMA of the previous block's end-state $\Sm$
(\S\ref{sec:sigma-swap}). This per-block freeze is the load-bearing
component of a novel defense against \emph{wedge attacks} --- mint
sandwiches and JIT minting that exploit the convexity of LMSR against
the linearity of proportional mint and burn (\S\ref{sec:wedge}). The
defense combines a raw single-block $\Delta\sigma_q$ deviation gate, a
min-clamped proportional burn that structurally closes JIT extraction, a
per-window rate limit on the proportional mint-resize factor $\gamma$
(the fraction by which a mint grows the pool, equal to the new LP minted
divided by the pre-mint supply), and a post-mint LP lock. Together
these allow LMSR pricing to coexist with the standard AMM treatment of
LP tokens as pro-rata claims on pool reserves. Pool parameters are immutable post-deployment; the
only mutable admin levers on a live pool are the protocol-fee recipient
and the irreversible \texttt{kill()} switch. The admin cannot mint LP,
change fees, freeze users, or move pool reserves --- a leaked admin key
presents little to no danger to LPs (\S\ref{sec:admin-risk}). Pool
creation is restricted to the trusted operator, who is responsible for
upstream token vetting via a documented validator pipeline. Where the
implementation makes sub-ulp rounding choices, every choice is biased
toward existing LPs and against takers and the protocol developers; the
protocol's primary client is the LP. This document specifies the pricing
kernel, the wedge-attack defense, the liquidity operations, the
production fee structure, the admin and deployment surface, and the
security-review methodology applied prior to launch.
\end{abstract}

\tableofcontents

\section{Overview}
\label{sec:overview}

Liquidity Party is a multi-asset AMM that holds many tokens in a single pool
and quotes every pair within that pool from a single convex potential. The
system is designed for capital efficiency on long-tail asset baskets: an
$n$-asset pool exposes $n(n-1)/2$ pairs with one set of LP shares,
eliminating the fragmentation cost of pairwise CFMM deployments. Per-operation
gas grows with the basket size (the LMSR potential touches every slot), so the
practical range is roughly $n \le 50$ tokens per pool; approximate gas costs
for production-typical sizes are summarized in Table~\ref{tab:gas}.

\begin{table}[h]
\centering
\caption{Approximate gas costs by pool size}
\label{tab:gas}
\begin{tabular}{rrrrrrr}
\toprule
Assets & Pairs & Swap & Mint & Burn & SwapMint & BurnSwap \\
\midrule
 2 &    1 & 138{,}000 &   174{,}000 &   143{,}000 &   173{,}000 &   144{,}000 \\
10 &   45 & 150{,}000 &   473{,}000 &   424{,}000 &   388{,}000 &   349{,}000 \\
20 &  190 & 165{,}000 &   847{,}000 &   776{,}000 &   655{,}000 &   608{,}000 \\
30 &  435 & 180{,}000 & 1{,}220{,}000 & 1{,}130{,}000 &   923{,}000 &   859{,}000 \\
50 & 1225 & 210{,}000 & 1{,}970{,}000 & 1{,}830{,}000 & 1{,}460{,}000 & 1{,}380{,}000 \\
\bottomrule
\end{tabular}
\end{table}

\paragraph{System layout.}
The on-chain system consists of four singleton-or-immutable contracts:

\begin{itemize}
  \item \texttt{PartyPlanner} --- factory and pool registry. \texttt{onlyOwner}
        \texttt{newPool} deploys each \texttt{PartyPool} via \texttt{CREATE2}
        with the caller-specified configuration vector: $(\kappa,
        \mathbf{f})$, the wedge-defense immutables (gate threshold
        \texttt{MINT\_DEVIATION\_PPM}, \texttt{EMA\_SHIFT\_BLOCKS}, the
        per-window $\gamma$ cap \texttt{MAX\_GAMMA\_PER\_WINDOW\_PPM}, and
        \texttt{MINT\_LOCK\_BLOCKS}), and the protocol-fee settings. The
        planner stores no default parameters of its own
        (\S\ref{sec:arch}). Indexes pools by token for off-chain discovery.
  \item \texttt{PartyPool} --- the pool itself, holding token reserves and
        issuing an ERC-20 LP share token. The pool is a non-upgradeable
        \texttt{DELEGATECALL} proxy over two implementation libraries
        (\texttt{PartyPoolExtraImpl1}, \texttt{PartyPoolExtraImpl2}) due to the
        EIP-170 24~KB contract-size cap.
  \item \texttt{PartyInfo} --- read-only view contract for prices, quotes, and
        denominators. Pricing helpers are exposed here rather than on the pool
        so that \texttt{PartyPool} fits its size budget.
  \item \texttt{PartyConcierge} --- user-facing router. Takes token addresses
        rather than indices, supports a single ERC-20 approval covering every
        pool, and exposes an EIP-7730 clear-signing surface. Funds pools via
        the callback funding mode (see \S\ref{sec:funding}), and offers an
        optional keeper-executed mint queue for rate-limited mints
        (\S\ref{sec:mint-queue}).
\end{itemize}

The pool's per-token configuration (bases and per-asset fee rates) is stored
in an SSTORE2-style data contract referenced by an immutable address in the
pool. Hot paths read it via \texttt{EXTCODECOPY}; no SLOAD on the bases or
fees occurs on the swap path.

\paragraph{Notation.}
We index assets by $i \in \{0,\dots,n-1\}$ and write the pool's internal
inventory as $\mathbf{q} = (q_0,\dots,q_{n-1})$, with the size metric
$\Sm = \sum_i q_i$. The liquidity parameter is $\bofq = \kappa\,\Sm$ for a
positive constant $\kappa$ fixed at deployment. Internal quantities are
maintained in Q64.64 fixed-point (ABDKMath64x64); user-facing token amounts
are integer base units (\texttt{uint256}) and are converted on the boundary
using per-asset denominators $\beta_i$ (the \emph{bases}) set at deployment:
$q_i = \mathrm{balance}_i / \beta_i$.

\paragraph{Primary audience: liquidity providers.}
Liquidity Party's first concern is the LP. The protocol is engineered so
that every decision the LP cannot make on their own behalf --- rounding,
approximation, cap-and-invert branch selection, fee/protocol-share split,
emergency response --- is biased in favor of the LP. Takers and integrators
are first-class users of the protocol surface but inherit the LP-favor
invariants from the LP side; they should design assuming that any
sub-ulp ambiguity in the rounding chain will resolve in the LP's favor
and against their own position.

\paragraph{Reading guide.}
\S\ref{sec:kernel} specifies the pricing kernel and its block-anchored
$b$. \S\ref{sec:wedge} specifies the wedge-attack defense (the
raw single-block $\Delta\sigma_q$ gate, the min-clamped burn, the
per-window $\gamma$ rate limit, and the post-mint LP lock), which is the novel
mechanism that lets LMSR pricing coexist with pro-rata LP accounting. \S\ref{sec:liquidity}
specifies the user-facing liquidity operations (proportional mint,
proportional burn, \texttt{swapMint}, \texttt{burnSwap}) as they
interact with the defense. \S\ref{sec:lp-protections} consolidates the
LP-favor guarantees that the kernel and the defense establish.
\S\ref{sec:fees}--\S\ref{sec:funding} describe fees and the four funding
modes. \S\ref{sec:arch} covers the deployed contract surface,
\S\ref{sec:admin} the admin powers and parameter fixity, and
\S\ref{sec:deploy} the operator's deployment-time calibration choices.
\S\ref{sec:security} documents the security review process that gates
production launch. \S\ref{sec:risk} states the residual risk surface and
operational mitigations, including the framing of token-slot drainage
as expected and industry-normal LP behavior.

\section{Pricing Kernel}
\label{sec:kernel}

\subsection{LMSR Cost Function and the $\kappa\Sm$ Parameterization}

The pool maintains the LMSR cost function in \emph{inventory convention}
\begin{equation}
  C(\mathbf{q};\,b) \;=\; -b\,\ln\!\left(\sum_{k=0}^{n-1} e^{-q_k/b}\right),
\end{equation}
where larger $q_i$ corresponds to a more abundant slot and therefore a
\emph{cheaper} marginal price. The liquidity parameter $b$ tracks pool
size as $b = \kappa\,\Sm$ but is held \emph{fixed for the duration of
each block} and anchored to the size metric computed at the end of the
previous block (\S\ref{sec:block-b}). This quasi-static LS-LMSR
treatment retains the three properties that drive the rest of the
design:

\begin{enumerate}
  \item \textbf{Scale invariance.} A proportional rescaling
        $\mathbf{q} \mapsto \lambda\mathbf{q}$ that also rescales $b$ by
        $\lambda$ leaves every pairwise marginal price unchanged, so
        proportional mint and burn are value-neutral at matched $b$.
  \item \textbf{Bounded LP loss.} The classical LMSR worst-case-loss
        $b\,\ln n$ becomes $\kappa\,\Sm\,\ln n$ at the current state; per
        unit of size this is $\kappa\,\ln n$, fixed by the deployment
        parameters.
  \item \textbf{Liquidity sensitivity.} As the pool absorbs trades and
        accrues fees across blocks, $\Sm$ grows and the next block's
        anchor grows with it; depth deepens automatically in proportion
        to TVL without governance intervention.
\end{enumerate}

The continuous LS-LMSR construction of Othman et al.\ (2013) admits a
similar size-dependent $b$ but with a continuously evolving liquidity
trajectory whose exact pricing requires integral evaluation. Our
$\bofq = \kappa\,\Sm$ choice is a fixed proportionality that is sampled
once per block from a block-end snapshot, eliminating the path-integral
dependence within any single transaction while retaining liquidity
sensitivity across blocks.

\subsection{Inventory Convention: The Convexity Flip vs Hanson's LMSR}
\label{sec:convention-flip}

Liquidity Party's pricing kernel is mathematically a sign-and-curvature
reflection of Hanson's original LMSR cost function. This subsection
states the relationship explicitly because it is easy to misread the
formula above as Hanson's cost function with a notational
substitution. It is not: it is a different cost surface that supports
a different agent.

\paragraph{Hanson's outstanding-shares convention.}
Hanson (2002) defines the LMSR cost function over a vector $\mathbf{s}
= (s_0, \dots, s_{n-1})$ of \emph{outstanding shares}, the cumulative
quantity of each outcome-share the market maker has sold into the
market:
\begin{equation}
  \label{eq:hanson-cost}
  C_H(\mathbf{s};\,b)
  \;=\;
  +\,b\,\ln\!\left(\sum_{k=0}^{n-1} e^{+s_k/b}\right).
\end{equation}
The cost function is \emph{convex} in $\mathbf{s}$. Its gradient
$\partial C_H / \partial s_i = e^{s_i/b} /
\sum_k e^{s_k/b}$ is the marginal price the market maker charges to
sell one additional share of outcome $i$: \emph{more abundant
outstanding shares of $i$ imply a higher marginal price for that
outcome}. The economic agent is a market-maker seat that is
\emph{short} the basket --- it sells shares into the market and
collects payments at the gradient price. Its worst-case running cost is
bounded by $b \cdot \ln n$. Hanson's LMSR is the canonical pricing
kernel for prediction markets, where the market-maker seat is paid for
information aggregation rather than for holding inventory.

\paragraph{Liquidity Party's inventory convention.}
A standard AMM does the opposite job. There is no market-maker seat;
the protocol holds a pool of $n$ asset reserves
$\mathbf{q} = (q_0, \dots, q_{n-1})$ and an LP-token holder is
\emph{long} that basket pro-rata. A trader pays asset $i$ \emph{into}
the pool (raising $q_i$) and pulls asset $j$ \emph{out} (lowering
$q_j$); the marginal price along that exchange should fall as $q_j$
shrinks and the pulled asset becomes scarce. We need a cost surface
under which $q_i$ is inventory, not outstanding shares, with the
correct ``more abundant slot is cheaper'' direction on the marginal
price.

The naive route --- substitute $s_i := -q_i$ into
\eqref{eq:hanson-cost} --- only flips the gradient direction; it leaves
the outer sign of $C_H$ positive and the function convex in
$\mathbf{q}$. The convexity-in-$\mathbf{q}$ would mean the cost
function \emph{grows} when the pool gets larger via a proportional
deposit, contradicting the desired LP-side accounting (a deposit
shouldn't read as a cost incurred). To get an inventory cost surface
whose gradient is correct \emph{and} whose curvature matches the
LP-as-long-the-basket geometry, we flip \emph{both} signs:
\begin{equation}
  \label{eq:inventory-cost}
  C(\mathbf{q};\,b)
  \;=\;
  -\,b\,\ln\!\left(\sum_{k=0}^{n-1} e^{-q_k/b}\right).
\end{equation}
This is \emph{not} the same function as \eqref{eq:hanson-cost} under
relabeling. The inner exponent is negated \emph{and} the outer sign is
flipped --- a double flip. The result is \emph{concave} in
$\mathbf{q}$, while Hanson's is convex in $\mathbf{s}$. The gradient
$\partial C / \partial q_i = e^{-q_i/b} / \sum_k e^{-q_k/b}$ is the
marginal weight assigned to slot $i$, with $\partial C / \partial q_i$
\emph{decreasing} in $q_i$ --- exactly the ``more abundant slot is
cheaper'' direction we want, and exactly the inverse of Hanson's
direction.

\paragraph{Two flips, not one.}
The pairing of \emph{input--output direction flip}
($s_i \leftrightarrow -q_i$, ``shares-sold'' becoming
``inventory-held'') with \emph{curvature flip} (convex $C_H$ becoming
concave $C$) is what distinguishes our formulation from a notational
relabel of Hanson:
\begin{itemize}
  \item The \textbf{input/output flip} reverses the semantic role of
        the state vector. In Hanson, $s_i$ rises when the AMM
        \emph{sells} into the market; in our convention, $q_i$ rises
        when a trader sells asset $i$ \emph{into} the pool. The same
        positive $\Delta$ on the state vector means opposite
        cash-flow direction on the protocol side.
  \item The \textbf{curvature flip} reverses the role of the cost
        function. Hanson's $C_H$ is the market-maker's cumulative
        \emph{liability} to share-holders, bounded above by
        $b\ln n$; our $C$ is a concave potential whose level sets
        define the cost-preserving swap manifold. The LP-side
        worst-case loss bound is still $b\ln n$, but it arrives from
        the opposite direction --- the maximum amount of value that
        can be \emph{drained} from a pro-rata-long pool by an
        adversarial price path, rather than the maximum payout the
        Hanson seat can owe.
\end{itemize}

\paragraph{Cost preservation under inventory convention.}
A swap of $a$ units of asset $i$ in for $y$ units of asset $j$ out
is value-neutral on the LP side if and only if it preserves the
inventory cost:
\begin{equation}
  \label{eq:cost-preserve}
  C(\mathbf{q} + a\,e_i - y\,e_j;\,b)
  \;=\;
  C(\mathbf{q};\,b).
\end{equation}
The closed-form solution at fixed $b$ is the Hanson exact-input
mapping rewritten in inventory variables:
\begin{equation}
  y \;=\; b\,\ln\!\Big(1 + r_0\,\big(1 - e^{-a/b}\big)\Big),
  \quad
  r_0 = e^{(q_j - q_i)/b}.
\end{equation}
Note the position of the minus signs: the marginal-price ratio $r_0$
is keyed to $(q_j - q_i)$ (output minus input), and the input
contribution is $e^{-a/b}$ (input depressing the output rate as it
floods slot $i$). These signs are the visible footprint of the double
flip in the cost surface; they would all run the other way under
Hanson's outstanding-shares convention.

\paragraph{The convexity flip under LS-LMSR.}
Othman et al.\ (2013) extend LMSR to liquidity-sensitive by letting
$b$ track a function of state, typically $b = \kappa \cdot \Sm$. Under
Hanson's convention this keeps $C_H$ convex; the larger $\Sm$ flattens
the convexity, deepening liquidity. Under our inventory convention
the same parameterization keeps $C$ concave, and the larger $\Sm$
flattens the concavity, deepening liquidity. The LS extension
respects the convention --- it does not introduce a third flip. The
liquidity parameter $b = \kappa \cdot \min(\sigma_{\mathrm{swap}},\,
\sigma_{\mathrm{live}})$ (\S\ref{sec:block-b}) carries the same
``larger $b$ is deeper liquidity'' meaning that Hanson and Othman et
al.\ attach to it, but it parameterizes a concave surface rather than
a convex one.

\paragraph{Why this matters operationally.}
The convexity flip is invisible to a trader who only reads the
marginal-price formula $P(i \to j) = \exp((q_j - q_i)/b)$. That
formula has the same shape as Hanson's $P_H(s_i)$, just with a different
sign on the state argument. The flip becomes visible whenever the LP
side of the accounting is exercised: proportional mint, proportional
burn, the LP-loss bound, the wedge-attack defense
(\S\ref{sec:wedge}), and the cost-preservation invariant the kernel
solves. Throughout this document, when we write ``LMSR cost
function,'' ``LS-LMSR,'' or ``cost-preservation,'' the quantity meant
is the inventory-convention $C$ of \eqref{eq:inventory-cost} --- not
Hanson's $C_H$.

\subsection{Marginal Prices}

With $b$ held constant within a block, the gradient of $C$ recovers
softmax-style prices and the instantaneous marginal exchange rate for
swapping input $i \to$ output $j$ is
\begin{equation}
  P(i \to j \mid \mathbf{q},\,b) \;=\; \exp\!\left(\frac{q_j - q_i}{b}\right).
\end{equation}
A more abundant output slot (larger $q_j$) buys more output tokens per
unit input. Per-pair denomination-adjusted prices are exposed via
\texttt{PartyInfo.price(pool,\,i,\,j)} as a Q128.128 fixed-point quantity
that accounts for the bases $\beta_i, \beta_j$.

\subsection{Block-Anchored $b$ via $\min(\sigma_{\mathrm{swap}}, \sigma_{\mathrm{live}})$}
\label{sec:block-b}

Every pricing read on the hot path --- the swap kernel, the per-leg
quotes inside \texttt{swapMint} and \texttt{burnSwap}, the
\texttt{PartyInfo} view helpers --- computes
\begin{equation}
  b \;=\; \kappa \cdot \min\!\big(\sigma_{\mathrm{swap}},\,\sigma_{\mathrm{live}}\big),
\end{equation}
where $\sigma_{\mathrm{live}} = \sum_i q_i$ is the live size metric
recomputed from the current pool inventory, and $\sigma_{\mathrm{swap}}$
is an EMA-anchored estimate of the previous block's end-state
$\sigma_{\mathrm{live}}$ (the full state machine and update rules are
specified in \S\ref{sec:sigma-swap}). The minimum ensures that any
mid-block expansion of $\sigma_{\mathrm{live}}$ --- whether from
arbitrage vig, proportional mint, or attacker skew --- does not
retroactively widen the depth the rest of the block prices against; a
block can only \emph{shrink} the effective $b$, never grow it.

\paragraph{Why fixed-$b$ within a block.}
Recomputing $b$ per leg of a multi-step operation (\texttt{swapMint},
\texttt{burnSwap}) from the evolving local state
$\mathbf{q}_{\mathrm{local}}$ would produce three mutually inconsistent
pricing surfaces across \texttt{swap}, \texttt{swapMint}, and
\texttt{burnSwap}, and closed cycles combining the three operations
could extract LP value across the inconsistency. Anchoring $b$ to the
same block-aligned $\min(\sigma_{\mathrm{swap}}, \sigma_{\mathrm{live}})$
for every leg of every operation collapses the three surfaces into one
and closes the cycle exploit; combined with the LP-favor rounding policy
(\S\ref{sec:rounding}), the residual asymmetry between the operations is
consistently in the LP's favor.

\paragraph{Why $\min$, not $\sigma_{\mathrm{swap}}$ directly.}
A pool that has been recently \emph{drained} by a proportional burn
sees $\sigma_{\mathrm{live}}$ drop below $\sigma_{\mathrm{swap}}$ (the
EMA still reflects the larger pre-burn state). Pricing against the
stale $\sigma_{\mathrm{swap}}$ would let an attacker swap into the
shrunken pool at the pre-burn depth --- buying at a liquidity depth the
pool no longer has, and extracting the difference from the remaining
LPs. The $\min$ closes that surface: post-burn the pool prices at the shrunken
$\sigma_{\mathrm{live}}$, and post-skew-swap it prices at the unmoved
$\sigma_{\mathrm{swap}}$ --- both directions key to the \emph{lower} of
the two, which is uniformly LP-favoring.

\subsection{Swap Mapping}
\label{sec:swap-mapping}

Let $i, j$ be the input and output asset indices and let $a$ be the
internal input passed to the kernel. The production swap kernel
(\texttt{LMSRKernel.\allowbreak swapAmountsForExactInput}) evaluates the
single-pass Hanson exact-input formula at the block-anchored $b$:
\begin{align}
  &b \;=\; \kappa \cdot \min\!\big(\sigma_{\mathrm{swap}},\,\sigma_{\mathrm{live}}\big), \nonumber\\
  &r_0 \;=\; \exp\!\big((q_j - q_i)/b\big), \nonumber\\
  &y \;=\; b\,\ln\!\Big(1 + r_0\,\big(1 - e^{-a/b}\big)\Big).
\end{align}
The implementation caps output at the available inventory $q_j$ when
the inner term of $\ln$ degenerates (cap-and-invert branch). There is
no path-integral correction inside the kernel: the block-fixed $b$
already provides the round-trip invariant the kernel needs, and any
intra-block closed sequence of \texttt{swap}, \texttt{swapMint}, and
\texttt{burnSwap} prices against a single consistent surface. Fees
(\S\ref{sec:fees}) and LP-favor rounding (\S\ref{sec:rounding}) make
every fee-aware closed cycle strictly loss-making for the attacker, and
the wedge defense (\S\ref{sec:wedge}) closes the
LMSR-convexity attack on mint and burn that the block-fixed surface
would otherwise leave open.

\subsection{LP-Safety Property of the Block-Anchored Kernel}
\label{sec:lp-safety}

The path-exact reference is the liquidity-sensitive LMSR mapping under
inventory convention: the $y^*$ that solves the cost-preservation
equation
\begin{equation}
  \label{eq:lslmsr}
  C\big(\mathbf{q} + a\,e_i - y^*\,e_j;\,b^*\big) \;=\; C(\mathbf{q};\,b^*),
\end{equation}
under a continuously evolving $b^* = \kappa\,\Sm$. With $b$ truly
state-dependent, both $b$ and the inner exponents change as the state
moves from $\mathbf{q}$ to $\mathbf{q} + a\,e_i - y\,e_j$, so $y^*$ has
no closed form. Our implementation freezes $b$ at
$\kappa\,\min(\sigma_{\mathrm{swap}}, \sigma_{\mathrm{live}})$ for the
duration of the block, with $\sigma_{\mathrm{live}}$ at most the
pre-trade size metric on the first state change of the block; subsequent
mid-block reads see the same or smaller $\sigma_{\mathrm{live}}$ value
(via the $\min$), so the effective $b$ used during a swap is at most the
exact-LS-LMSR midpoint trajectory $b$.

\paragraph{LP-safety property.}
\emph{For every $(a, \mathbf{q}, \kappa, n)$ in the production envelope,
the kernel output $y$ satisfies}
\[
  y \;\le\; y^*.
\]
\emph{Equivalently, the swapper receives at most what the path-exact
LS-LMSR would pay, and the LP loses at most what the path-exact LS-LMSR
would surrender.} The proof is the inequality on $b$ above: the
block-fixed $b$ never exceeds the path-exact $b$-trajectory, the Hanson
exact-in mapping is increasing in $b$ on the relevant regime, and
therefore the kernel's $y$ is below the path-exact $y^*$.

\paragraph{Empirical zero-extraction.}
The complementary guarantee --- that no closed \emph{sequence} of legs
(e.g.\ \texttt{swap}, \texttt{swapMint}, \texttt{burnSwap},
\texttt{mint}, \texttt{burn}) can accumulate sub-ulp residue into an
extractable round trip --- is addressed by the adversarial cycle test
suite in \texttt{test/SwapMintBurnSwapBAnchorExploit.t.sol} and the skew-attack
suites (\texttt{ArbHarness}, \texttt{StandardPoolAttacks},
\texttt{PoC\_WedgeAttack}). The suites run closed cycles against pools
already driven to imbalance, covering 2-leg pure-swap loops, 3- and
4-leg mixed swap / \texttt{swapMint} / \texttt{burnSwap} /
\texttt{mint} / \texttt{burn} loops, and intra-block self-sandwich
cycles. The assertion at every cycle boundary is \emph{strict}: any
closed cycle that leaves the attacker even \textbf{1 wei} richer
(measured at pre-cycle marginal prices) is a test failure, and the same
strict bound is required to hold on the pool TVL at pre-cycle prices.
\textbf{No grid point fails.} The shared-block-$b$ surface makes the
intra-block cycle pricing self-consistent; the
wedge defense (\S\ref{sec:wedge}) closes the
mint-side convexity attack on top of that.

\paragraph{Capacity caps.}
For output $y$ that would exceed available $q_j$, the kernel returns
$y = q_j$ at the cap-and-invert boundary. The inverse mapping
(\texttt{amountInForExactOutput}) returns the input required for a
target output, using
\begin{equation}
  a(y) \;=\; b \,\ln\!\left(\frac{r_0}{r_0 + 1 - e^{y/b}}\right),
\end{equation}
which is the exact closed-form inverse of the Hanson exact-in mapping
(used internally on the cap-and-invert branch and on the inverse leg of
\texttt{swapMint}, not as a public exact-out swap entry point).

\paragraph{Numerical stability.}
The kernel computes ratios like $r_0 = \exp((q_j - q_i)/b)$ directly
rather than dividing exponentials, guards every \texttt{exp} and
\texttt{ln} call to bounded domains (\texttt{EXP\_LIMIT = 32.0}), and
checks the positivity of the inner $\ln$ argument before evaluation.
Domain violations revert; outputs are clamped to capacity rather than
extrapolated. The cost function evaluates log-sum-exp with one-pass
dynamic recentering for stability under dispersed $q_k$.

\subsection{Slippage Protection}

Every user-facing swap accepts a \texttt{minAmountOut} (or, on
\texttt{swapMint}, a \texttt{maxAmountIn}). If the computed output is below
\texttt{minAmountOut}, the call reverts. This is denomination-aware slippage
protection requiring no price-unit conversions. Off-chain callers that prefer
a price-based bound may use the bisection helper
\texttt{PartyInfo.swapAmountsForExactPrice} to convert a target post-swap
forward price into a \texttt{(maxAmountIn, minAmountOut, fee)} triple suitable
for direct submission to \texttt{swap()}.

\section{Wedge-Attack Defense}
\label{sec:wedge}

LMSR pricing is convex; proportional mint and burn are linear in pool
size. The mismatch between the two creates a structural attack family
we call the \emph{wedge}: an attacker front-runs a victim's mint with a
skew swap that moves the pool to an off-market state, the victim's
proportional deposit deepens the pool at the manipulated price ratio,
the attacker back-runs with a reverse swap that captures the
convexity-bonus value that the deposit added to the wedge. The
attacker extracts face value from both the victim and the pre-existing
LPs.

This attack is the central obstacle to running LMSR pricing alongside
a standard AMM-style LP token, where LP shares are pro-rata claims on
pool reserves. Hanson's original LMSR is operated by a market-maker
seat, not by pro-rata LPs; the convexity belongs to the seat, not to
depositors, so the wedge has no mint/burn surface to attack. Liquidity
Party retains LMSR's marginal-price geometry while issuing pro-rata LP
tokens, which re-opens the wedge surface. The defense specified in
this section is what allows that combination.

A full design and parameter justification is in
\texttt{doc/rate-limited-mints.md}; this section gives the
production-load-bearing summary.

\subsection{Defense Architecture}
\label{sec:wedge-arch}

Four components, plus a load-bearing coupling to the pool's swap fee:

\begin{enumerate}
  \item \textbf{The raw single-block $\Delta\sigma_q$ gate.} A single
        Q64.64 slot, \texttt{\_prevBlockEndSigmaQ}, holds the
        $\sigma_q = \Sm$ snapshot captured at the first state change of
        the current block --- the end-of-previous-block value, by
        construction. The gate trips and \texttt{mint} /
        \texttt{swapMint} revert while
        \[
          \big|\sigma_{\mathrm{live}} - \sigma_{\mathrm{prevBlockEnd}}\big|
          \;\ge\;
          \tau \cdot \sigma_{\mathrm{prevBlockEnd}},
        \]
        where $\tau$ is the immutable per-pool gate threshold (the
        \texttt{MINT\_DEVIATION\_PPM} parameter, $\tau_d$). The
        reference is a fixed block-\emph{start} snapshot, not in-block
        state, so spam-fill within a block cannot manipulate it --- all
        in-block activity is measured against the same anchor. Crucially
        the gate measures only the move \emph{within this block}: it
        trips inside a block that itself carries a $\ge \tau_d$ jump and
        reopens the next, rather than staying tripped while a slow EMA
        catches up (\S\ref{sec:gate}, ``Why raw, not $\sigma_{\mathrm{swap}}$'').
  \item \textbf{Min-clamped proportional burn.} The proportional-burn
        payout is scaled by
        $\alpha' = \alpha \cdot \min(\sigma_{\mathrm{swap}},\,
        \sigma_{\mathrm{live}}) / \sigma_{\mathrm{live}}$ rather than
        pure $\alpha$. When $\sigma_{\mathrm{swap}}$ lags
        $\sigma_{\mathrm{live}}$ --- as it does immediately after a
        swap added vig --- burners receive their \emph{pre-swap} pool
        value, not the post-swap value. The vig stays in the pool for
        LPs who wait for $\sigma_{\mathrm{swap}}$ to converge. This is
        what closes JIT minting structurally.
  \item \textbf{Per-window $\gamma$ rate limit.} A continuous,
        EMA-decayed accumulator caps $\sum \gamma$ across all mints
        within the current EMA window at
        \texttt{MAX\_GAMMA\_PER\_WINDOW\_PPM}. This bounds the
        aggregate gate-boundary residual per window and prevents
        attacker self-sandwich cycling.
  \item \textbf{Post-mint LP lock.} Freshly-minted LP is non-burnable
        for a per-pool immutable window of \texttt{MINT\_LOCK\_BLOCKS}
        (\S\ref{sec:mint-lock}). The LP remains freely transferable, but
        the lock travels with the tokens --- a transfer of still-locked
        LP migrates the lock to the recipient, who cannot burn it until
        the original unlock block. The lock forces a mint to commit
        capital across blocks before it can exit, which closes the
        atomic \texttt{mint}$\to$\texttt{burn} round-trip that would
        otherwise exhaust the rate-limit budget at zero net capital
        cost, and also closes multi-block JIT around large swaps once
        $\sigma_{\mathrm{swap}}$ has caught up enough that the min-clamp
        no longer penalizes the burn.
\end{enumerate}

The fourth load-bearing piece is the pool's per-asset \textbf{swap
fee} (\S\ref{sec:fees}). A linear-in-swap-size fee on the attacker's
two skew-and-reverse legs flips the gate-permitted boundary-attack
per-attack net negative; the per-family $\tau$ and fee
recommendations in \texttt{doc/rate-limited-mints.md} are calibrated
together against the boundary-attack simulator.

\subsection{State and EMA Update Rules}
\label{sec:sigma-swap}

The defense state lives in three storage slots adjacent to the LMSR
state in \texttt{PartyPoolStorage.PoolState}:

\begin{itemize}
  \item \texttt{\_sigmaSwap} (Q64.64): the EMA-smoothed estimate of
        $\sigma_q$. It is \emph{not} the gate reference --- it drives
        only the block-anchored swap $b$ (\S\ref{sec:block-b}) and the
        burn value clamp (\S\ref{sec:min-gated-burn}), where its slow
        convergence is a feature.
  \item \texttt{\_prevBlockEndSigmaQ} (Q64.64): the $\sigma_q$ snapshot
        captured at the first state-change of each new block (which is
        the end-of-previous-block snapshot, by construction). \textbf{This
        is the mint-gate reference} --- both the deviation anchor and the
        denominator (\S\ref{sec:gate}). It also serves as the EMA target
        for the $\sigma_{\mathrm{swap}}$ step.
  \item \texttt{\_lastUpdateBlock} (uint64): the block number of the
        last $\sigma_{\mathrm{swap}}$ step.
  \item \texttt{\_gammaAccum} (Q64.64): the EMA-decayed accumulator of
        applied $\gamma$ for the per-window rate limit.
  \item \texttt{\_gammaAccumLastBlock} (uint64): block number of the
        last accumulator decay update.
\end{itemize}

The $\sigma_{\mathrm{swap}}$ state evolves by the following rules:

\paragraph{On the first state change of a new block.}
The pool first \emph{captures} the gate reference
$\sigma_{\mathrm{prevBlockEnd}} = \sigma_{\mathrm{live}}$ (read from
\texttt{qInternal} before any mutation in this call, so it is exactly
the end-of-previous-block value), then steps the EMA toward it:
\[
  \sigma_{\mathrm{swap}} \;\mathrel{+}=\;
  \frac{\sigma_{\mathrm{prevBlockEnd}} - \sigma_{\mathrm{swap}}}
       {2^{\text{\texttt{EMA\_SHIFT\_BLOCKS}}}}.
\]
The freshly-captured $\sigma_{\mathrm{prevBlockEnd}}$ is the fixed
anchor against which every mint in this block is gated, so a plain mint
with no prior same-block activity sees $\sigma_{\mathrm{live}} =
\sigma_{\mathrm{prevBlockEnd}}$ and passes the gate by construction.
The shift is fixed at one shift per block-with-activity, regardless of
how many blocks have elapsed since the last update. Quiet pools
accumulate no ``convergence credit''. After many quiet blocks the
$\sigma_{\mathrm{swap}}$ defense is in fact \emph{stronger}, not
weaker.

\paragraph{On proportional mint.}
$\sigma_{\mathrm{swap}}$, $\sigma_{\mathrm{live}}$, \emph{and} the gate
reference $\sigma_{\mathrm{prevBlockEnd}}$ all scale by $(1 + \gamma)$.
Scaling $\sigma_{\mathrm{swap}}$ preserves the proportional invariant
$\sigma_{\mathrm{swap}} / \sigma_{\mathrm{live}}$ for its b-anchor and
burn-clamp roles. Scaling the gate reference by the \emph{same} factor
is what keeps the raw gate from reading a proportional event as
volatility: a proportional grow preserves relative prices, so without
this write a second same-block mint would see the first mint's
$(1+\gamma)$ inventory scaling as a $\Delta\sigma_q$ jump and spuriously
trip. The scaling is multiplicative (not an overwrite with
$\sigma_{\mathrm{live}}$), so any genuine same-block skew already in the
reference scales along with it and stays caught --- the same-block
sandwich defense is preserved.

\paragraph{On proportional burn.}
The pool computes the effective burn fraction $\alpha'$ as defined in
\S\ref{sec:wedge-arch}, then scales reserves $\mathbf{q}$ by
$(1 - \alpha')$ (the value actually paid out) while scaling \emph{both}
$\sigma_{\mathrm{swap}}$ and the gate reference
$\sigma_{\mathrm{prevBlockEnd}}$ by $(1 - \alpha)$ --- the
\emph{LP-supply} scale, not the reserve scale. The LP-token supply is
likewise reduced by the full $\alpha$ that the user requested. Tying the
$\sigma_{\mathrm{swap}}$ shrink to $(1 - \alpha)$ rather than
$(1 - \alpha')$ is what makes a clamped burn
\emph{fragmentation-invariant}: $N$ burns of $\Delta L/N$ total exactly
the same payout as one burn of $\Delta L$
(\S\ref{sec:min-gated-burn}). The gate reference scales by the same
$(1 - \alpha)$ for the analogous reason as on mint: a burn shrinks
inventory proportionally, so a later same-block mint must not read that
shrink as a $\Delta\sigma_q$ jump. When the clamp is inactive ($\alpha'
= \alpha$ --- the converged, full-drain, or killed cases) the reserve
and supply scale factors coincide.

\paragraph{On swap.}
Only $\sigma_{\mathrm{live}}$ changes (it is derived from
$\mathbf{q}$ on demand). No mutation to $\sigma_{\mathrm{swap}}$
happens within a block; the block's end-state $\sigma_{\mathrm{live}}$
is what feeds the next block's EMA step.

\paragraph{Non-volatility reconciliation.}
Two events move $\sigma_{\mathrm{live}}$ without being a manipulation
jump, and every LP operation reconciles them into the gate reference so
a same-block mint does not read them as in-block volatility.
\emph{Retained LP-fee backlog}: a plain \texttt{swap} advances
\texttt{qInternal} by the gross output but retains the LP-fee share
off-inventory, so the backlog must be folded back into \texttt{qInternal}
before the next LP operation gates against it. \texttt{burn},
\texttt{swapMint}, and \texttt{burnSwap} fold it explicitly at their
start; \texttt{mint} folds it implicitly, by rescaling
$\sigma_{\mathrm{swap}}$ across the proportional rebuild (which absorbs
the backlog) rather than by the bare $(1 + \gamma)$ factor. When the
fold is explicit, $\sigma_{\mathrm{prevBlockEnd}}$ is advanced
\emph{additively} by the same $\sigma_{\mathrm{live}}$ delta --- additive,
not an overwrite, so any genuine prior same-block swap stays measured
against the true block-start reference. \emph{Physical donations /
over-delivery drift} ($\mathrm{balance} > \mathrm{cached} + \mathrm{owed}$)
are folded by the \texttt{mint} / \texttt{burn} sweep, rescaling
$\sigma_{\mathrm{swap}}$ by the resulting ratio (the hot \texttt{swap}
and the single-asset variants skip the sweep for gas; their drift is
reclaimed by the next \texttt{mint} / \texttt{burn}). Outside these
reconciliations, the block-start capture, and the proportional mint/burn
scale, $\sigma_{\mathrm{prevBlockEnd}}$ is never written.

\paragraph{Anti-spam-fill, anti-quiet-pool.}
The step is keyed on \texttt{block.number}, not on the count of
in-block transactions or wall-clock time. Spam-filling 256 swaps in a
block produces exactly one EMA step. After arbitrary inactivity, the
first state change in a new block applies one shift step against the
end-of-previous-block snapshot --- the EMA does not ``cash in''
accumulated convergence on a single update. The
\texttt{block.number} anchor also rules out L2-sequencer
\texttt{block.timestamp} manipulation.

\subsection{The Gate}
\label{sec:gate}

\texttt{mint} and \texttt{swapMint} check the gate at the start of the
operation, comparing the \emph{current} (post-swap-leg, for
\texttt{swapMint}) $\sigma_{\mathrm{live}}$ against the block-start
reference $\sigma_{\mathrm{prevBlockEnd}}$:
\[
  \big|\sigma_{\mathrm{live}} - \sigma_{\mathrm{prevBlockEnd}}\big| \cdot 10^6
  \;\ge\;
  \texttt{MINT\_DEVIATION\_PPM} \cdot \sigma_{\mathrm{prevBlockEnd}}
  \;\Rightarrow\;
  \text{revert \texttt{"volatile market"}.}
\]
$\sigma_{\mathrm{prevBlockEnd}}$ is used as \emph{both} the deviation
anchor and the denominator, and the $\sigma_{\mathrm{live}}$ here is the
\emph{pre-mint} $\sigma_q$: the mint's own proportional $\gamma$ growth
is deliberately excluded, so the gate asks ``did the pool move this
block before this mint'', not ``does this mint grow the pool''.

The inequality is \emph{non-strict} (`$\ge$', not `$>$') for an
important reason: an attacker can size a \texttt{swapMint}'s swap leg
to land the pool exactly at $\tau \cdot \sigma_{\mathrm{prevBlockEnd}}$
deviation from the block-start reference. A strict `$>$' check would
let that exactly-on-threshold poison swap through; non-strict catches
the boundary case.

Gate-passing is required but not sufficient: each \texttt{mint} also
checks \texttt{maxAmountsIn} and the rate-limit budget; the operation
reverts if either fails.

\paragraph{Why raw, not $\sigma_{\mathrm{swap}}$.}
Anchoring this gate on the EMA-smoothed $\sigma_{\mathrm{swap}}$ instead
($|\sigma_{\mathrm{live}} - \sigma_{\mathrm{swap}}| \ge \tau \cdot
\sigma_{\mathrm{swap}}$) would conflate two unrelated things: slow
organic LVR drift accumulated across many blocks, and single-block
manipulation. After \emph{any}
genuine $\sigma_q$ move, $\sigma_{\mathrm{live}}$ jumps immediately but
$\sigma_{\mathrm{swap}}$ only catches up at $2^{-\texttt{EMA\_SHIFT\_BLOCKS}}$
per block, so the deviation stays above $\tau$ for $\approx
2^{\texttt{EMA\_SHIFT\_BLOCKS}}$ blocks and the gate stays
\emph{tripped that whole time}, locking out honest mints long after the
move was over (on a realistic 12\,s jump tail this tripped honest mints
$\sim$72\% of the time with a $\sim$98-minute worst-case lockout). The
raw single-block signal separates the two: organic $\sigma_q$ is flat
between arbitrage jumps, so $|\sigma_{\mathrm{live}} -
\sigma_{\mathrm{prevBlockEnd}}|$ is $\approx 0$ in a quiet block
regardless of how far the pool has drifted over the day. The gate
trips only \emph{in} a block carrying a $\ge \tau_d$ jump and reopens
the next --- organic trip rate $\approx$ 3.8\% (volatile family) to
$\approx$ 0.03\% (correlated-peg family), worst-case lockout 1--4
blocks. $\sigma_{\mathrm{swap}}$ is retained for the jobs its slow
convergence is good at: the lagging $b$-anchor and the burn clamp.

\subsection{Min-Clamped Burn and Why It Closes JIT}
\label{sec:min-gated-burn}

Burn is always available --- the gate never disables it --- but the
\emph{amount returned} is min-clamped. Concretely, for a requested burn
fraction $\alpha = \mathit{lpAmount} / \mathit{lpSupply}$:
\begin{itemize}
  \item \textbf{Full drain} (\texttt{lpAmount == lpSupply}):
        $\alpha' = \alpha = 1$, the burner receives the entire pool
        inventory. The min-clamp is bypassed because there is no
        remaining LP to receive the gated-out value.
  \item \textbf{Killed pool} (\texttt{\_killed} is set): $\alpha' =
        \alpha$, pure proportional withdrawal. The min-clamp would
        unfairly penalize LPs caught with pending vig at kill time;
        the value would never reach honest LPs anyway since further
        pool activity is frozen.
  \item \textbf{Normal case:}
        $\alpha' = \alpha \cdot \min(\sigma_{\mathrm{swap}},\,
        \sigma_{\mathrm{live}}) / \sigma_{\mathrm{live}}$. The user
        receives $\alpha'$ of every reserve token, while the LP-token
        burn is at the full $\alpha$.
\end{itemize}

The pool inventory $\mathbf{q}$ scales by $(1 - \alpha')$ --- the
fraction actually paid out --- while \emph{both} the LP-token supply
and $\sigma_{\mathrm{swap}}$ scale by $(1 - \alpha)$. Scaling
$\sigma_{\mathrm{swap}}$ by the LP-supply factor $(1 - \alpha)$ rather
than the reserve factor $(1 - \alpha')$ is deliberate, and the next
paragraph shows why.

\paragraph{Why $(1-\alpha)$: fragmentation invariance.}
With the clamp active ($\sigma_{\mathrm{swap}} < \sigma_{\mathrm{live}}$),
a naive $(1 - \alpha')$ scaling of $\sigma_{\mathrm{swap}}$ would
preserve the $\sigma_{\mathrm{swap}}/\sigma_{\mathrm{live}}$ ratio
across the burn, but it would let a fast LP \emph{fragment} their exit
into many small burns and recapture part of the convergence-window
penalty: each split chunk would face the same clamp ratio against a
proportionally richer per-LP pool, so the sum of $N$ small payouts
would exceed the single-burn payout. Scaling $\sigma_{\mathrm{swap}}$ by
$(1 - \alpha)$ instead holds the value-per-LP quantity
\[
  x \;=\; \frac{\sigma_{\mathrm{swap}} \cdot R}{\sigma_{\mathrm{live}} \cdot S}
\]
invariant across the burn ($R$ the reserve value, $S$ the LP supply):
the two numerator factors scale by $(1 - \alpha)$ and the two
denominator factors by $(1 - \alpha')$, so every chunk in a fragmented
exit --- and any third party burning in the same block --- sees exactly
the $x$ of the pre-burn pool. The aggregate of $N$ split burns equals
$\Delta L \cdot x = \alpha' \cdot R$, the single-burn payout. The
mechanism is a \emph{tightening} of the clamp on each subsequent
same-block burn: because $\sigma_{\mathrm{swap}}$ shrinks more than
$\sigma_{\mathrm{live}}$, the next burner faces a stricter ratio, which
is exactly the lever that keeps the missed vig with the slower-burning
LPs. When $\alpha' = \alpha$ (no clamp, full drain, or killed) the
$(1-\alpha)$ and $(1-\alpha')$ factors coincide and this reduces to
plain proportional scaling.

\paragraph{Why JIT is closed structurally.}
Consider the attack: mint $\gamma$, observe an incoming swap that
adds vig $V$ to $\sigma_q$, burn $\alpha$ in the same block.
\begin{itemize}
  \item Pre-attack: $\sigma_q = \sigma_{q,\mathrm{pre}}$,
        $\sigma_{\mathrm{swap}} = \sigma_{q,\mathrm{pre}}$ (converged).
  \item Post-mint: $\sigma_{\mathrm{swap}} = (1 + \gamma)\,
        \sigma_{q,\mathrm{pre}} = \sigma_{\mathrm{live}}$ (the mint
        scales both proportionally).
  \item Post-swap: $\sigma_{\mathrm{live}} = (1 + \gamma)\,
        \sigma_{q,\mathrm{pre}} + V$;
        $\sigma_{\mathrm{swap}}$ is unchanged (no mid-block EMA step).
  \item Burn $\alpha = \gamma / (1 + \gamma)$:
        $\alpha' = \alpha \cdot \sigma_{\mathrm{swap}} /
        \sigma_{\mathrm{live}}$;
        the payout is
        $\alpha' \cdot \sigma_{\mathrm{live}} = \alpha \cdot
        \sigma_{\mathrm{swap}} = \gamma \cdot \sigma_{q,\mathrm{pre}}$.
\end{itemize}
The attacker receives back exactly their mint deposit, with no share
of the vig $V$. The $V$ stays in pool inventory and accrues to LPs who
do \emph{not} burn within the convergence window. No protocol mint fee
is needed for this defense; honest LP entry incurs no protocol fee.

\paragraph{What honest LPs see.}
LPs that hold longer than roughly $4 \cdot 2^{\texttt{EMA\_SHIFT\_BLOCKS}}$
blocks ($\approx$ 3--14 hours depending on shift) see $\alpha' \approx
\alpha$. There is no exit penalty in the steady state. LPs forced to
exit immediately after a vig-adding swap pay the min-clamp as an
effective exit fee that decays over the EMA window, with the
``missed'' vig staying for slower-burning LPs.

\subsection{Per-Window Rate Limit}
\label{sec:rate-limit}

The accumulator \texttt{\_gammaAccum} decays continuously per block:
\[
  \texttt{\_gammaAccum}
  \;\mathrel{\mathord{*}{=}}\;
  \left(1 - 2^{-\texttt{EMA\_SHIFT\_BLOCKS}}\right)^{\!\Delta},
  \quad \Delta = \text{blocks since last update,}
\]
implemented as $\Delta$ linear shifts (each shift subtracts
$\texttt{\_gammaAccum} \gg \texttt{EMA\_SHIFT\_BLOCKS}$ from itself).
This matches the $\sigma_{\mathrm{swap}}$ window scale, so the
boundary-attack residual budget and the gate's natural recovery
share a single timescale.

Each \texttt{mint} or \texttt{swapMint}:
\begin{enumerate}
  \item Decays \texttt{\_gammaAccum} up to the current block.
  \item Computes
        \(\gamma_{\mathrm{remaining}} =
          \texttt{MAX\_GAMMA\_PER\_WINDOW\_PPM} - \texttt{\_gammaAccum}\).
  \item Sets
        \(\gamma_{\mathrm{fill}} =
          \min(\gamma_{\mathrm{request}},\,\gamma_{\mathrm{remaining}})\).
  \item Reverts if $\gamma_{\mathrm{fill}} \le 0$ (no budget left).
  \item Reverts with \texttt{"rate limited"} if
        $\gamma_{\mathrm{fill}} < \gamma_{\mathrm{request}}$ and the
        caller passed \texttt{partialFillAllowed = false}; otherwise
        proceeds with the partial fill.
  \item After the mint, adds $\gamma_{\mathrm{fill}}$ to
        \texttt{\_gammaAccum}.
\end{enumerate}

Per-window rather than per-mint because a per-mint cap is cosmetic
under transaction splitting; the per-window cap binds
$\sum \gamma_i$ regardless of how many mints there are or who submits
them. With JIT closed structurally by the min-clamped burn,
\texttt{MAX\_GAMMA\_PER\_WINDOW\_PPM} is a backstop against
attacker self-sandwich cycling rather than the primary JIT defense; the
standard pools set it to $4{,}000$~PPM ($0.4\%$ of supply per window) on
the volatile OG pool and $10{,}000$~PPM ($1\%$) on the correlated-peg
pool, which comfortably accommodates organic LP growth without
re-opening the cycling surface.

\subsection{Post-Mint LP Lock}
\label{sec:mint-lock}

The rate limit caps the \emph{rate} of LP issuance, but a per-window cap
alone is exhaustible at zero net cost: an attacker can atomically
\texttt{mint} $\gamma$ and \texttt{burn} it back in the same block,
keeping the rate-limit credit (the accumulator already booked
$\gamma_{\mathrm{fill}}$) while the min-clamped burn returns the deposit
at par (at the moment of attack $\sigma_{\mathrm{swap}} =
\sigma_{\mathrm{live}}$, so $\alpha' = \alpha$). On an idle pool this
collapses to a gas-only griefing DoS on honest minters. The post-mint
LP lock closes this surface.

\paragraph{Mechanism.} Every \texttt{mint} and \texttt{swapMint}
appends a time-locked \emph{cohort} to the receiver: the minted LP
becomes non-burnable until block
$\texttt{block.number} + \texttt{MINT\_LOCK\_BLOCKS}$. The lock is
enforced in the ERC-20 \texttt{\_update} hook --- a \emph{burn} that
would dip into the receiver's still-locked balance reverts, while a
\emph{transfer} of still-locked LP is permitted but carries the lock to
the recipient (see \emph{Cohort accounting} below), so the unlock block
is preserved no matter who holds the tokens. The
window \texttt{MINT\_LOCK\_BLOCKS} is a per-pool immutable, capped at
$50{,}400$ blocks (one week of L1 blocks) at the planner; the suggested
L1 default is $300$ blocks ($\approx$ 1 hour), scaled up proportionally
on shorter-block L2s to preserve the wall-clock deterrent.
\texttt{initialMint} does not lock --- it is operator-controlled and no
attacker action precedes it.

\paragraph{Cohort accounting and transfers.} Each account holds a
FIFO list of cohorts, kept sorted by unlock block and lazily pruned as
they expire; the list is bounded at \texttt{MAX\_LOCK\_ENTRIES} $= 32$
live cohorts per account (a new mint to an account already at the cap
reverts until a cohort expires). The lock travels with the tokens: a
transfer that moves still-locked LP migrates the minimum FIFO prefix of
cohorts to the recipient, each retaining its original unlock block, so
the lock cannot be shed by forwarding LP to a fresh address. Subsequent
mints never extend an existing cohort --- each mint locks only its own
freshly-issued LP.

\paragraph{Deterrence.} Any non-zero window closes the same-block
round-trip outright. Deterrence against \emph{sustained} DoS scales
with the window: the integrated locked capital an attacker must keep
deposited to hold the rate limit pinned at $\Gamma_{\max}$ grows
roughly linearly with the lock duration (at \texttt{EMA\_SHIFT\_BLOCKS}
$=4$, a one-hour lock pins on the order of $9\,\Gamma_{\max}$ worth of
$\gamma$ continuously deposited). Honest LPs that mint to hold are
unaffected, as are aggregators and zappers that mint LP for an end user
to keep --- the lock only bites a position that tries to exit within
the window.

\subsection{Boundary Attack and the Swap-Fee Coupling}

The gate caps an attacker's per-bundle skew at $\delta_f \le \tau
\cdot \sigma_q$. Through the LMSR mint-sandwich mechanism, this lets
them extract at most $\tau \cdot \gamma/(1+\gamma) \cdot \sigma_q$ per
attack at zero fee. With the pool's per-asset swap fee applied to
both the forward and reverse skew legs, the per-attack net flips
negative: the swap fees on the two legs exceed the gate-bounded gross
extraction at the per-family $(\tau, f)$ recommendations in
\texttt{doc/rate-limited-mints.md} (and verified by the
boundary-attack simulator). The boundary attack is therefore
\emph{closed} (net-negative for the attacker), not merely bounded.

\subsection{Synchronous Semantics and the Optional Mint Queue}

\texttt{mint} and \texttt{swapMint} are synchronous at the pool: the
user submits one transaction and either receives LP in the same
transaction, sees a partial fill (if they opted in), or sees the call
revert with one of the operational reverts (\texttt{"volatile market"},
\texttt{"rate limited"}, \texttt{"slippage control"},
\texttt{"insufficient funds"}). The pool holds no queue and runs no
deferred execution itself. Because the rate limit can force a partial
fill, the \texttt{PartyConcierge} router offers an \emph{optional} mint
queue with permissionless keeper execution as a convenience
(\S\ref{sec:concierge}); a user who opts out sees the plain synchronous
semantics above.

\subsection{Attack Surfaces Closed and Residual Surfaces}

\paragraph{Closed.}
\begin{enumerate}
  \item Mint sandwich (gate caps skew + swap fee makes the per-attack
        residual net-negative).
  \item Spam-fill EMA bypass (block-anchored step ignores intra-block
        transaction count).
  \item Quiet-pool single-shot convergence (fixed one-shift-per-block
        step regardless of elapsed blocks).
  \item JIT mint (min-clamped burn structurally returns the attacker's
        deposit without any vig share).
  \item Burn-side sandwich (verified arithmetically self-griefing for
        any attacker LP share $\lambda_a \le 1 - \alpha$, see
        \texttt{doc/rate-limited-mints.md} ``Burn safety'').
  \item Atomic \texttt{mint}$\to$\texttt{burn} rate-limit exhaustion
        (the post-mint LP lock forces every mint to commit capital
        across blocks before it can exit, \S\ref{sec:mint-lock}).
  \item Burn fragmentation (the $(1-\alpha)$ scaling of
        $\sigma_{\mathrm{swap}}$ makes a clamped exit
        fragmentation-invariant, so splitting a burn cannot recapture
        the convergence-window penalty, \S\ref{sec:min-gated-burn}).
\end{enumerate}

\paragraph{Residual.}
\begin{enumerate}
  \item \textbf{Sustained DOS griefing.} An adversary willing to burn
        capital indefinitely can keep the gate tripped by continuously
        re-skewing. Per the threat model this is unprofitable
        extraction for the attacker; honest mints simply revert and
        retry. No tokens get stuck.
  \item \textbf{LP-as-oracle.} The gate controls mints; pair prices
        can still swing within $\exp(\tau)$ under a gate-permitted
        in-bounds skew. At the standard-pool threshold $\tau =
        30$~PPM (\texttt{MINT\_DEVIATION\_PPM} on both the OG and
        correlated-peg pools) this is approximately $1.00003\times$,
        negligible. Downstream consumers that read pool marginal price
        for liquidations should still apply their own TWAP.
  \item \textbf{Initial-mint window.} The very first mint occurs in
        the same transaction as \texttt{initialMint}, which is
        operator-controlled; no attacker action is possible before
        initial mint.
  \item \textbf{Rate-limit DOS.} An attacker can exhaust
        \texttt{MAX\_GAMMA\_PER\_WINDOW\_PPM} by minting their own LP
        to grief honest users. The post-mint LP lock
        (\S\ref{sec:mint-lock}) removes the zero-cost atomic version of
        this: held pressure on the rate limit now requires capital
        locked for \texttt{MINT\_LOCK\_BLOCKS}, with the integrated
        locked capital growing roughly linearly in the lock window.
        What remains is a capital-bound DOS, not a value-extraction
        surface.
\end{enumerate}

\section{Liquidity Operations}
\label{sec:liquidity}

The pool issues an ERC-20 LP share token. The total LP supply $L$ tracks
$\Sm$ proportionally: the deployer specifies the initial supply $L^{(0)}$
as an argument to \texttt{initialMint} (defaulting to $10^{18}$ if zero is
passed), independent of pool size, and subsequent operations preserve the
ratio $L / \Sm$ on every proportional change. The LP-share ERC-20 reports
\texttt{decimals() = 18} regardless of $L^{(0)}$; the deployer typically
chooses $L^{(0)}$ to land the pool at a convenient initial share price
relative to the initial deposit basket.

\subsection{Proportional Mint}
\label{sec:mint}

The \texttt{mint} entry point is gated by the wedge defense
(\S\ref{sec:wedge}) --- the raw $\Delta\sigma_q$ gate and the
per-window rate limit --- and partial-fill aware. Its full signature
is
\begin{quote}\small
\texttt{mint(payer, fundingSelector, receiver, lpTokenAmount,
maxAmountsIn, minLpOut, partialFillAllowed, deadline, cbData)
\hfill\(\to\) (lpMinted, gammaFilled).}
\end{quote}
The caller specifies a target LP amount $\Delta L =$
\texttt{lpTokenAmount}. The pool:

\begin{enumerate}
  \item Captures the gate reference $\sigma_{\mathrm{prevBlockEnd}}$
        and steps $\sigma_{\mathrm{swap}}$ if this is the first state
        change of a new block (\S\ref{sec:sigma-swap}).
  \item Decays the per-window accumulator
        \texttt{\_gammaAccum} (\S\ref{sec:rate-limit}).
  \item Checks the raw $\Delta\sigma_q$ gate, comparing the pre-mint
        $\sigma_{\mathrm{live}}$ against the block-start reference
        $\sigma_{\mathrm{prevBlockEnd}}$ (\S\ref{sec:gate}); reverts
        \texttt{"volatile market"} if tripped.
  \item Computes the requested
        $\gamma_{\mathrm{request}} = \Delta L / L$, clamps to
        $\gamma_{\mathrm{fill}} = \min(\gamma_{\mathrm{request}},\,
        \gamma_{\mathrm{remaining}})$, reverts \texttt{"rate limited"}
        if $\gamma_{\mathrm{fill}} \le 0$ or if $\gamma_{\mathrm{fill}}
        < \gamma_{\mathrm{request}}$ and the caller did not pass
        \texttt{partialFillAllowed = true}.
  \item For each token $j$, computes the deposit needed as
        $\gamma_{\mathrm{fill}} \cdot q_j$, reverts
        \texttt{"slippage control"} if any deposit exceeds the
        per-token \texttt{maxAmountsIn[j]} bound.
  \item Pulls funds via the chosen funding mode
        (\S\ref{sec:funding}). Any over-delivery from a funding
        callback is absorbed into reserves (donated to LPs) per the
        rounding policy (\S\ref{sec:rounding}).
  \item Applies the proportional mint: $q_j \mapsto (1 +
        \gamma_{\mathrm{fill}})\,q_j$ for each $j$, and mints
        $\gamma_{\mathrm{fill}} \cdot L$ LP shares to \texttt{receiver}.
        Reverts \texttt{"slippage control"} if
        $\mathit{lpMinted} < \mathit{minLpOut}$.
  \item Scales $\sigma_{\mathrm{swap}}$ \emph{and} the gate reference
        $\sigma_{\mathrm{prevBlockEnd}}$ proportionally (both
        $\mathrel{*}= (1 + \gamma_{\mathrm{fill}})$,
        \S\ref{sec:sigma-swap}) and credits the accumulator
        ($\texttt{\_gammaAccum} \mathrel{+}= \gamma_{\mathrm{fill}}$).
  \item Appends a post-mint lock cohort on the minted LP at
        \texttt{receiver}, non-burnable for \texttt{MINT\_LOCK\_BLOCKS}
        (the LP stays transferable; the lock migrates with it,
        \S\ref{sec:mint-lock}).
  \item Returns $(\mathit{lpMinted},\, \mathit{gammaFilled})$. For
        typical mints inside the window budget, $\mathit{lpMinted} =
        \Delta L$.
\end{enumerate}

There is \textbf{no protocol mint fee}; honest LP entry is free. The
JIT defense lives entirely on the burn side
(\S\ref{sec:min-gated-burn}).

\subsection{Proportional Burn}
\label{sec:burn}

\texttt{burn(payer, receiver, lpAmount, minAmountsOut, deadline,
unwrap)} burns $\Delta L =$ \texttt{lpAmount} LP shares and returns the
per-token withdraw amounts. The full payout rule, including the
JIT-defense min-clamp, is specified in \S\ref{sec:min-gated-burn}; the
summary is:
\begin{itemize}
  \item The pool computes the requested burn fraction
        $\alpha = \Delta L / L$ and the effective payout fraction
        $\alpha'$ per \S\ref{sec:min-gated-burn} (with the full-drain
        and killed-pool edge cases bypassing the min-clamp).
  \item The pool sends $\alpha' \cdot q_k \cdot \beta_k$ uint of each
        token $k$ to \texttt{receiver}; if \texttt{unwrap} is set and
        token $k$ is the configured native wrapper, the output is
        unwrapped to native currency before delivery.
  \item Per-token \texttt{minAmountsOut[k]} bounds are enforced;
        revert string \texttt{"slippage control"}.
  \item The reserves $\mathbf{q}$ scale by $(1 - \alpha')$ (the
        amount paid out), while the LP-token supply and
        $\sigma_{\mathrm{swap}}$ both scale by $(1 - \alpha)$. The
        $(1-\alpha)$ scaling of $\sigma_{\mathrm{swap}}$ makes a clamped
        burn fragmentation-invariant (\S\ref{sec:min-gated-burn}).
\end{itemize}

\texttt{burn} is \emph{not} gated by the deviation check and is
\emph{not} disabled by \texttt{kill()}; LPs can always exit pro rata.

Proportional mint and burn are value-neutral in the fee-free kernel at
matched $b$: every pairwise marginal price is invariant under
$\mathbf{q} \mapsto \lambda\mathbf{q}$ when the corresponding $b$ also
scales by $\lambda$.

\subsection{Single-Asset Mint: \texttt{swapMint}}
\label{sec:swapmint}

\texttt{swapMint} is exact-LP-out: the caller specifies a target
$\Delta L$ and a \texttt{maxAmountIn} cap on the deposit of a single
input asset $i$. The operation is equivalent to executing the
externally observable sequence ``swap $i \to j$ for every $j \ne i$,
then proportional mint with the resulting basket.'' The fee-free
kernel computes the required input in one forward pass against a
memory-only local state $\mathbf{q}_{\mathrm{local}}$, with $b$
anchored to the same block-aligned $\kappa \cdot
\min(\sigma_{\mathrm{swap}},\,\sigma_{\mathrm{live}})$ used by the
\texttt{swap} kernel (\S\ref{sec:block-b}):

\begin{enumerate}
  \item Compute $\gamma = \Delta L / L$ and
        $\beta = \gamma/(1+\gamma) \in (0,1)$.
  \item Initialize $\mathbf{q}_{\mathrm{local}} = \mathbf{q}$ and
        anchor $b = \kappa \cdot \min(\sigma_{\mathrm{swap}},\,
        \sigma_{\mathrm{live}})$ once for the entire chain.
  \item For each $j \ne i$ in fixed index order, evaluate
        $r_{0,j} = \exp((q_{\mathrm{local}}^j -
        q_{\mathrm{local}}^i)/b)$ and solve the Hanson exact-out
        inverse for the cost of acquiring $\beta\,q_j$ of asset $j$:
        \[
          x_j \;=\; b\,\ln\!\left(
            \frac{r_{0,j}}{r_{0,j} + 1 - e^{\beta q_j / b}}\right).
        \]
        Update $\mathbf{q}_{\mathrm{local}}$:
        $q_{\mathrm{local}}^i \mathrel{+}{=} x_j$,
        $q_{\mathrm{local}}^j \mathrel{-}{=} \beta\,q_j$.
  \item Return the total deposit
        $\displaystyle a \;=\; \frac{\beta\,q_i + \sum_{j \ne i} x_j}{1 - \beta}.$
\end{enumerate}

The chain mutates a memory copy only; the pool's storage is updated
at the end with the new reserves and the minted LP shares. Crucially,
\emph{$b$ does not re-evaluate per leg}: every leg of the multi-step
chain prices against the same block-aligned $b$ as the standalone
\texttt{swap} kernel, collapsing the three operations
(\texttt{swap}, \texttt{swapMint}, \texttt{burnSwap}) onto a single
intra-block pricing surface and closing the cycle-extraction surface
that per-leg re-anchoring would otherwise open
(\S\ref{sec:block-b}, ``Why fixed-$b$ within a block'').

\paragraph{End-to-end flow.}
The pool runs the raw $\Delta\sigma_q$ gate, the rate-limit decay, and the
funding-mode pull around the kernel call, with the same partial-fill
and revert semantics as \texttt{mint}:
\begin{enumerate}
  \item Step $\sigma_{\mathrm{swap}}$ and capture
        $\sigma_{\mathrm{prevBlockEnd}}$ if first-of-block; absorb any
        retained fee backlog into the gate reference
        (\S\ref{sec:sigma-swap}); decay \texttt{\_gammaAccum}.
  \item Pull the input token via the chosen funding mode; apply the
        full swap-leg fee --- the input asset's fee \emph{plus} the
        other assets' fees, the same combined rate a decomposed
        multi-leg \texttt{swap} would charge (\S\ref{sec:fees}) --- and
        send the fee-discounted amount through the kernel.
  \item Apply the rate-limit cap to $\gamma_{\mathrm{request}}$;
        revert \texttt{"rate limited"} if $\gamma_{\mathrm{fill}} <
        \gamma_{\mathrm{request}}$ and \texttt{partialFillAllowed =
        false}.
  \item Check the raw $\Delta\sigma_q$ gate against the post-swap-leg
        $\sigma_{\mathrm{live}}$ relative to
        $\sigma_{\mathrm{prevBlockEnd}}$ --- the non-strict $\ge$
        catches the \emph{poison-swapMint at the gate boundary}
        described in \S\ref{sec:gate}; revert \texttt{"volatile market"}
        if tripped.
  \item Apply the proportional mint; update $\sigma_{\mathrm{swap}}$
        and \texttt{\_gammaAccum}; append a post-mint lock cohort on the
        minted LP (\S\ref{sec:mint-lock}); revert
        \texttt{"slippage control"} if \texttt{maxAmountIn} or
        \texttt{minLpOut} fails; emit \texttt{SwapMint}.
\end{enumerate}
Failure of any step reverts the entire operation; tokens remain with
the caller (or never leave them, in the Permit2 and callback cases).

\paragraph{Rounding.}
The exact-out cost $x_j$ uses ABDK floor primitives that round in the
trader's favor under the natural
\texttt{mul}/\texttt{div}/\texttt{exp}/\texttt{ln} semantics. To
restore the LP-favoring invariant, the kernel applies explicit
ceiling helpers (\texttt{\_ceilMul}, \texttt{\_ceilDiv},
\texttt{\_ceilExp}, \texttt{\_ceilLn}) at every step whose floor
direction would shave a sub-ulp gift to the swapper, and adds a 1-ulp
cushion in the asymptotic $\mathrm{EXP\_LIMIT}$ branch. Every
individual rounding decision is documented in-source by signed
partial derivative against the LP-favoring direction.

\subsection{Single-Asset Redeem: \texttt{burnSwap}}
\label{sec:burnswap}

\texttt{burnSwap} is the inverse of \texttt{swapMint}: the caller
burns $\Delta L$ LP shares and receives a single output asset $j$. The
operation is \emph{not} subject to the $\Delta\sigma_q$ gate
or the rate limit --- both apply to mint-side pool expansion, not to
withdrawal --- but the proportional-burn portion of the payout is
subject to the same min-clamp as standalone \texttt{burn}
(\S\ref{sec:min-gated-burn}). The fee-free kernel computes the payout
by simulating a proportional withdraw followed by swaps of every
withdrawn $a_k$ back into asset $j$ against the already-burnt local
state, with $b$ anchored to the block-aligned $\kappa \cdot
\min(\sigma_{\mathrm{swap}},\,\sigma_{\mathrm{live}})$ for every leg:

\begin{enumerate}
  \item Compute $\alpha = \Delta L / L$ and the effective
        $\alpha' = \alpha \cdot \min(\sigma_{\mathrm{swap}},\,
        \sigma_{\mathrm{live}}) / \sigma_{\mathrm{live}}$
        (with edge cases per \S\ref{sec:min-gated-burn}).
  \item Anchor $b = \kappa \cdot \min(\sigma_{\mathrm{swap}},\,
        \sigma_{\mathrm{live}})$ once for the entire chain.
        Initialize $\mathbf{q}_{\mathrm{local}} = (1 -
        \alpha')\,\mathbf{q}$ and start the payout with the
        proportional share $Y_j = \alpha'\,q_j$.
  \item For each $k \ne j$: set $a_k = \alpha'\,q_k$ (the withdrawn
        portion of asset $k$). Evaluate
        $r_{0,k} = \exp((q_{\mathrm{local}}^j -
        q_{\mathrm{local}}^k)/b)$ and the Hanson exact-in payout
        \[
          y_{k \to j} \;=\; b\,\ln\!\Big(1 + r_{0,k}\,\big(1 -
          e^{-a_k/b}\big)\Big).
        \]
        Accumulate $Y_j \mathrel{+}{=} y_{k \to j}$ and update
        $\mathbf{q}_{\mathrm{local}}$. If $y_{k \to j}$ exceeds
        available $q_{\mathrm{local}}^j$, cap the payout to
        $q_{\mathrm{local}}^j$ and derive the actually-consumed
        input via the exact-out inverse.
  \item Per-asset numerical degeneracy (non-positive inner $\ln$
        argument, $\mathrm{EXP\_LIMIT}$ overshoot) is handled by the
        asymptotic branch in
        \texttt{LMSRKernel.swapAmountsForBurn} described in
        \S\ref{sec:kernel}, contributing the bounded asymptotic
        payout rather than aborting the redeem.
\end{enumerate}

As with \texttt{swapMint}, \emph{$b$ does not re-evaluate per leg}.
The standalone \texttt{swap}, \texttt{swapMint}, and \texttt{burnSwap}
operations all see the same block-aligned $b$, so intra-block closed
cycles compose against one consistent pricing surface. The full
swap-leg fee --- the output asset's fee \emph{plus} the other assets'
fees --- is applied to the gross output (\S\ref{sec:fees}).

\paragraph{Kill-resistance asymmetry.}
\texttt{burnSwap} is disabled by \texttt{kill()}; \texttt{burn} is
not. A killed pool therefore permits proportional redemption but not
single-asset redemption. \texttt{burnSwap} also pays the full swap-leg
fee on its swap-back leg; the proportional \texttt{burn} does not.

\section{LP Protection Mechanisms}
\label{sec:lp-protections}

The pricing kernel (\S\ref{sec:kernel}) and liquidity operations
(\S\ref{sec:liquidity}) carry per-primitive rounding and approximation
choices that resolve every sub-ulp ambiguity in the LP's favor. This
section consolidates those guarantees and adds the admin-side LP
protections that constrain the operator surface. Each mechanism is
enforced in source; this is an index, not a statement of intent.

\subsection{Kernel-Level LP Safety}

\paragraph{Block-anchored $b$ never overshoots LS-LMSR.}
The production swap kernel
(\texttt{LMSRKernel.\allowbreak swapAmountsForExactInput}) evaluates
the single-pass Hanson exact-input formula at $b = \kappa \cdot
\min(\sigma_{\mathrm{swap}},\,\sigma_{\mathrm{live}})$
(\S\ref{sec:block-b}, \S\ref{sec:lp-safety}). The output paid to the
swapper is bounded above by the output the path-exact LS-LMSR mapping
would pay under the same inputs: the block-fixed $b$ is at most the
path-exact $b$-trajectory, and the Hanson exact-in mapping is
monotone increasing in $b$, so the kernel's output is at or below the
path-exact reference. The swapper \emph{never} receives more than the
path-exact LS-LMSR; the LP \emph{never} loses more than the path-exact
LS-LMSR.

\paragraph{Single intra-block pricing surface.}
Standalone \texttt{swap}, the per-leg sub-swaps inside
\texttt{swapMint}, and the per-leg sub-swaps inside \texttt{burnSwap}
all evaluate at the same block-aligned $b$. Closed intra-block cycles
combining the three operations therefore price against one consistent
surface and cannot extract value across pricing inconsistencies. The
deployed test suites \texttt{test/SwapMintBurnSwapBAnchorExploit.t.sol},
\texttt{ArbHarness}, \texttt{StandardPoolAttacks}, and
\texttt{PoC\_WedgeAttack} verify the strict-zero-extraction property
across the recommended operating envelope.

\paragraph{Capacity caps.}
Output is hard-capped at available inventory $y \le q_j$ in the
kernel. The cap-and-invert branch derives the input that the LP
actually absorbs, so the pool is never required to over-pay relative
to its inventory.

\paragraph{Single-asset chain LP-favor rounding.}
The \texttt{swapMint} and \texttt{burnSwap} chains
(\S\ref{sec:swapmint}--\S\ref{sec:burnswap}) carry a per-primitive rounding
discipline. Where ABDK's default \texttt{mul}/\texttt{div}/\texttt{exp}/\texttt{ln}
floor would shave a sub-ulp gift to the swapper, the kernel substitutes a
ceiling helper (\texttt{\_ceilMul}, \texttt{\_ceilDiv}, \texttt{\_ceilExp},
\texttt{\_ceilLn}); where the floor direction already serves LP-favor, the
default is kept. Every individual decision is documented in-source by signed
partial derivative against the LP-favoring direction. The
\texttt{EXP\_LIMIT} branch in \texttt{swapMint} carries a 1-ulp cushion on
top of the ceiling chain so that future ABDK-level rounding changes cannot
re-open a leak.

\subsection{Fee and Protocol-Share Rounding}

\begin{itemize}
  \item \textbf{Total fees round up} for the pool and against the taker.
  \item \textbf{Protocol share rounds down} for the LP and against the
        protocol developers.
  \item \textbf{Outputs to users floor}; \textbf{inputs from users ceil}.
  \item \textbf{Pre-funded dust} is credited to LP reserves rather than
        refunded to the caller.
\end{itemize}

The protocol developers absorb the rounding loss on the protocol-share
split; the LP keeps the residual. This is deliberate. Section
\ref{sec:fees} specifies the fee mathematics this rounding policy applies
to.

\subsection{Correlated-Asset LVR Reduction (Conditional)}

A structural LP benefit of the multi-asset LMSR geometry is a closed-form
\emph{cross-correlation discount} on the loss-versus-rebalancing (LVR)
rate (Olson \& Claude Opus, 2026, in-tree
\texttt{doc/lvr-multiasset-lmsr.pdf}): the more correlated the pool's
basket, the more the LMSR's shared-simplex coupling cancels the
relative-mode arbitrage that bleeds value from separable bilateral
CFMMs. For perfectly correlated equal-volatility baskets the
relative-mode contribution vanishes entirely.

This protection is \emph{conditional}: it is realized through pool
composition, not enforced by the protocol on every pool, and a pool of
uncorrelated assets does not receive it. The operator-side mechanics ---
basket-selection criteria, the discount formula, and quantitative
comparison of candidate compositions --- are in \S\ref{sec:token-selection}.

\subsection{Admin-Side LP Protection}
\label{sec:admin-risk}

The admin (planner owner / pool owner) cannot reach LP reserves. The
exhaustive CAN/CANNOT inventory is in \texttt{doc/security/admin-powers.md};
the LP-relevant summary is:

\begin{itemize}
  \item Admin \textbf{cannot} mint LP out of band, burn user LP, change
        $\kappa$, change fees on a live pool, pause user-side
        \texttt{burn}, blacklist or freeze user balances, upgrade the pool
        implementation, renounce ownership, or move funds out of the pool.
  \item Admin \textbf{can} deploy new pools, set the protocol-fee recipient
        on a live pool, and invoke the irreversible \texttt{kill()} switch
        on a pool they own.
\end{itemize}

\paragraph{Compromised-key worst case.}
A leaked admin key cannot reach LP reserves. Worst case: the attacker
calls \texttt{kill()} (disabling new deposits, swaps, and single-asset
withdrawals) and rotates \texttt{protocolFeeAddress} so any
not-yet-collected protocol fees route to the attacker on the next sweep.
Both \texttt{burn} (proportional LP redemption) and
\texttt{collectProtocolFees} (protocol-fee sweep) remain operational on a
killed pool, so LPs can always withdraw their pro-rata share of pool
reserves and any party may push already-accrued protocol fees to the
configured \texttt{protocolFeeAddress}. The core team retains the ability
to \texttt{kill()} the system as a final defensive measure --- even if the
admin key is compromised, the system can be put into the LP-can-only-exit
state. \textbf{A leaked admin key presents little to no danger to LP
reserves}; it is a denial-of-service surface and a partial protocol-fee
redirect surface, not a fund-loss surface.

The admin key is held on a hardware wallet (see
\S\ref{sec:deployments}).

\subsection{Burn Is Always Live}

\texttt{burn} (the proportional-redemption path) is not gated by
\texttt{kill()}. The proportional-redemption guarantee is preserved
across every kill scenario, every admin-key scenario, and every
operator-misconfiguration scenario by construction: there is no code path
in the protocol that can disable \texttt{burn} for an LP holding the LP
token. \texttt{collectProtocolFees} is likewise un-killable, so
already-accrued protocol fees are also recoverable on a killed pool ---
they sweep to the configured \texttt{protocolFeeAddress} regardless of
kill state. These are the only two state-mutating entry points that
remain live after \texttt{kill()}; every other mutating call reverts.

The one restriction on \texttt{burn} is the post-mint LP lock
(\S\ref{sec:mint-lock}): LP is non-burnable for
\texttt{MINT\_LOCK\_BLOCKS} after it is minted. This is a temporary,
self-incurred, time-bounded restriction on freshly-minted shares --- not
an admin lever and not a kill-state effect --- and it expires
automatically. Any LP held past its lock window, including all LP on a
killed pool once the window elapses, redeems pro rata without
qualification.

\section{Fee Structure}
\label{sec:fees}

\paragraph{Direction summary.}
Each user-facing entry point charges its fee on exactly one side of the
operation. \texttt{swap} and \texttt{burnSwap} are
\textbf{fee-on-output}: the fee-free kernel computes the gross payout
first, then the fee is deducted from that payout before the user
receives it. \texttt{swapMint} is \textbf{fee-on-input}: the fee is
deducted from the user-submitted input before the fee-free kernel
runs, so the kernel prices the net (post-fee) input. The proportional
\texttt{mint} and proportional \texttt{burn} entry points are
\textbf{fee-free}: no swap fee applies to either side, and JIT
extraction is closed structurally by the min-clamped burn
(\S\ref{sec:min-gated-burn}) rather than by a fee. The collected fee
is retained in the pool, in the token on the fee'd side, and split
between the protocol share (\S\ref{sec:protocol-fees}) and the LP
share (kept in reserves, raising $\Sm$ and the LP-share-denominated
value).

\subsection{Per-Asset Swap Fees}

Each asset $k$ has an immutable fee rate $f_k$ set at deployment,
expressed in parts-per-million. The deployer enforces $f_k <
10{,}000$~ppm ($< 1\%$). For an asset-to-asset \texttt{swap} $i \to
j$, the effective pair fee is the additive composition
\begin{equation}
  f_{\mathrm{eff}}(i, j) \;=\; f_i + f_j,
\end{equation}
applied to the gross kernel output on the output side:
\[
  y_{\mathrm{net}} \;=\; y_{\mathrm{gross}}\,(1 - f_{\mathrm{eff}}).
\]
The full user-submitted input $a_{\mathrm{user}}$ is forwarded to the
fee-free kernel, which returns the gross output $y_{\mathrm{gross}}$;
the user receives exactly $y_{\mathrm{net}}$ in token $j$, and the pool
retains $y_{\mathrm{gross}} - y_{\mathrm{net}}$ in token $j$ as the
collected fee (then split into protocol/LP shares).

The additive form is preferred over the exact multiplicative
$1 - (1 - f_i)(1 - f_j)$ because the cross-term at the deployment cap
($f_i, f_j < 10^{-2}$) is below $10^{-4}$ of the pair fee, and the
simpler formula is cheaper to evaluate and easier for users to reason
about.

\subsection{Compound-Operation Fees on \texttt{swapMint} and \texttt{burnSwap}}

\texttt{swapMint} and \texttt{burnSwap} each carry a swap leg (the
named asset is traded against all the others) and a fee-free
proportional leg. The swap leg is charged the \emph{full} swap-leg
fee, not just the named asset's own rate. Writing $f_m$ for the named
asset's fee (the input $i$ on \texttt{swapMint}, the output $j$ on
\texttt{burnSwap}) and $n$ for the basket size, the combined rate is
\begin{equation}
  f_{\mathrm{leg}}
  \;=\; f_m \;+\;
  \Big\lceil \tfrac{1}{n-1}\textstyle\sum_{k \ne m} f_k \Big\rceil,
\end{equation}
the named asset's fee plus the (ceiling-rounded) average of the other
assets' fees. This matches what the equivalent decomposed multi-leg
\texttt{swap} path --- the named asset traded into each of the $n-1$
others --- would charge under the additive pair-fee composition
$f_i + f_j$ of \S\ref{sec:fees}. For a two-asset pool it reduces
exactly to $f_i + f_j$. The fee \emph{side} matches the side on which
the user moves tokens:

\begin{itemize}
  \item \texttt{swapMint} is \textbf{fee-on-input}: the pool computes
        $a_{\mathrm{eff}} = a_{\mathrm{user}}\,(1 - f_{\mathrm{leg}})$
        and passes $a_{\mathrm{eff}}$ to the fee-free kernel, which
        mints exactly $\Delta L$ LP shares (exact-LP-out) for the
        gate-and-rate-limited fraction $\gamma_{\mathrm{fill}}$. The
        fee $a_{\mathrm{user}} - a_{\mathrm{eff}}$ is collected in the
        deposited token.
  \item \texttt{burnSwap} is \textbf{fee-on-output}: the fee-free
        kernel computes the gross payout $y_{\mathrm{gross}}$ for the
        burned LP fraction (post the min-clamp scaling,
        \S\ref{sec:min-gated-burn}); the user receives
        $y_{\mathrm{net}} = y_{\mathrm{gross}}\,(1 - f_{\mathrm{leg}})$.
        The fee $y_{\mathrm{gross}} - y_{\mathrm{net}}$ is retained in
        the withdrawn token.
\end{itemize}

In both cases the protocol share is split off and the remainder stays
in pool reserves. There is no separate LP-token-side fee on either
operation, and no protocol fee on the proportional \texttt{mint} or
\texttt{burn}. Charging the \emph{full} swap-leg fee --- rather than
only the named asset's rate --- closes a fee-bypass loophole: otherwise
an attacker could pair a compound operation with the fee-free
proportional \texttt{burn} or \texttt{mint} and pay only one side's
per-asset fee.

\subsection{Protocol Fees}
\label{sec:protocol-fees}

A protocol-fee share $\phi$ (ppm) is taken from every collected fee. The
deployer enforces $\phi < 300{,}000$~ppm ($< 30\%$). Accrued protocol fees
are tracked in a per-token ledger
(\texttt{\_protocol\allowbreak FeesOwed}) excluded from the pool's
liquidity-bearing balance (\texttt{\_cached\allowbreak UintBalances});
pending protocol fees do not participate in pricing or earn further
fees. The \texttt{collectProtocolFees()} entry point sweeps the ledger
to the configured \texttt{protocolFeeAddress} and zeroes it. The
recipient is read at call time, so a recipient rotation applied between
accrual and collection retroactively redirects pending fees (Uniswap
V3-style semantics); operators that need atomic handoff must collect
before rotating.

\paragraph{Public, unguarded sweep.}
\texttt{collectProtocolFees()} is intentionally \emph{permissionless} and
carries no signature, owner-gate, or other authorization check beyond the
reentrancy guard. Funds always sweep to the admin-set
\texttt{protocolFeeAddress} regardless of which address called the
function. This is a deliberate operational simplification: the sweep can
be triggered by any keeper, integrator, or monitoring bot without
coordinating a signature from the operator's hardware-wallet key, while
the destination remains under admin control via the standard
\texttt{setProtocolFeeAddress} surface. The only address that benefits
from the sweep is the address the admin chose; no caller can divert
proceeds.

\subsection{Rounding Policy}
\label{sec:rounding}

All rounding decisions are biased to protect existing LP value:

\begin{itemize}
  \item Total fees are rounded \emph{up} (taker pays the extra ulp; LP keeps
        it).
  \item The protocol share is rounded \emph{down} (LP keeps the residual; the
        protocol developers absorb the ulp).
  \item Pool outputs to users are floored; pool inputs from users are ceiled.
  \item Pre-funded amounts above the exact required amount are credited to
        reserves rather than refunded (LP gift).
  \item Q64.64 ceil helpers (\texttt{\_ceilMul}, \texttt{\_ceilDiv},
        \texttt{\_ceilExp}, \texttt{\_ceilLn}) are applied per-primitive in
        \texttt{swapMint}/\texttt{burnSwap} to restore LP favor where ABDK's
        default floor would otherwise leak a sub-ulp gift to the swapper.
\end{itemize}

\section{Funding Modes}
\label{sec:funding}

The user-facing entry points that pull external tokens into the
pool --- \texttt{swap}, \texttt{mint}, and \texttt{swapMint} --- accept any
of four token-delivery modes via a \texttt{fundingSelector} parameter.
The kernel mathematics are identical across modes.

\paragraph{No funding mode on burn paths.}
\texttt{burn} and \texttt{burnSwap} do \emph{not} take a
\texttt{fundingSelector} and do not pull external tokens. The asset the
caller spends on these operations is the LP token itself, and the pool
\emph{is} the LP token's issuing contract --- the LP-token debit is a
storage write in the same contract that holds the pool reserves, so no
ERC-20 \texttt{transferFrom}, callback, or Permit2 signature is required
to move value into the pool on a burn. The caller's LP balance (or an
existing LP-token allowance) is the only authorization surface. This is
why none of the four funding modes below applies to either burn entry
point.

\begin{description}
  \item[\textsf{APPROVAL}] Standard ERC-20 \texttt{transferFrom} from the
        caller. Caller must hold the input balance and have approved the pool.
        The mode requires \texttt{msg.sender == payer} to prevent
        third-party-allowance abuse.
  \item[\textsf{PREFUNDING}] Caller deposits tokens to the pool address before
        invoking the entry point; the pool measures the balance delta against
        \texttt{\_cachedUintBalances + \_protocolFeesOwed}. Excess above the
        required amount is gifted to LPs. Native ETH is auto-wrapped via the
        configured wrapper when the pool's input token equals the wrapper
        address and the pool's ETH balance covers the amount. This delta is
        \emph{unauthenticated}: it carries no depositor identity, and the
        \texttt{msg.sender == payer} gate only prevents abuse of a victim's
        \emph{allowance}, not consumption of pre-deposited tokens. PREFUNDING is
        therefore intended only for integrating contracts that bundle the
        transfer and the entry-point call \emph{atomically in one transaction};
        a non-atomic deposit (transfer in one tx, call in another) can be
        front-run and consumed by any caller who sets \texttt{payer = msg.sender}.
        This residual race is a deliberate design choice, not a defended
        invariant --- EOA and cross-transaction integrations should use
        \textsf{APPROVAL}, \textsf{PERMIT2}, or callback funding instead.
  \item[\textsf{CALLBACK}] Aggregators and routers implement a callback
        interface. The pool calls
        \texttt{\textit{fundingSelector}(nonce, token, amount, cbData)}
        on the \texttt{payer} address and measures the resulting balance
        delta to determine the funded amount.
  \item[\textsf{PERMIT2}] Uniswap Permit2 signature-based transfer with
        witness data binding the operation's parameters to the signature.
        Witness types are defined per entry point in
        \texttt{PartyPoolPermit2Witness} for \texttt{swap}, \texttt{mint},
        and \texttt{swapMint}.
\end{description}

\paragraph{Native ETH.} The pool accepts native ETH on the entry points
declared \texttt{payable}, but only the configured wrapper (a WETH-style
contract) may push raw ETH back into the pool via \texttt{receive()}. Any
other direct ETH transfer reverts so stray sends cannot be stranded. Multi-
asset mint loops carry an explicit per-call ETH budget so the wrap branch
cannot over-claim against a constant \texttt{msg.value}.

\section{Contract Architecture}
\label{sec:arch}

\subsection{PartyPool: Proxy and Storage}

\texttt{PartyPool} is a non-upgradeable \texttt{DELEGATECALL} proxy.
Two implementation libraries (\texttt{PartyPoolExtraImpl1},
\texttt{PartyPoolExtraImpl2}) host the body of \texttt{initialMint},
\texttt{mint}, \texttt{burn}, \texttt{swapMint}, \texttt{burnSwap},
\texttt{collectProtocolFees}, and the deployment-time
\texttt{init} routine. The pool's own runtime is a thin facade so its
deployed bytecode plus its creation-code embed in
\texttt{PartyPoolInitCode} both fit under EIP-170's 24~KB
contract-size cap. Implementation addresses are immutable; there is no
upgrade path.

Storage layout is mirrored by \texttt{PartyPoolStorage.PoolState} for use
under \texttt{DELEGATECALL}. Per-token \texttt{bases} and \texttt{fees} are
not held in storage --- they live in an SSTORE2 ``BFStore'' data contract
deployed at \texttt{init} time, whose address is captured in the
\texttt{IMMUTABLE\_BFSTORE} pool immutable. Hot paths read it via
\texttt{EXTCODECOPY}, avoiding bounds-check SLOADs on the array length.

The LMSR state \texttt{(\_lmsr.kappa, \_lmsr.qInternal)} is a cache: the
pool re-derives the post-swap, post-mint, and post-burn \texttt{qInternal}
from the cached uint balances and the immutable bases on every mutation,
keeping the kernel state synchronized with physical reserves across all
state-changing paths.

\subsection{PartyPlanner: Factory and Registry}

\texttt{PartyPlanner.newPool} (\texttt{onlyOwner}) deploys a new
\texttt{PartyPool} via \texttt{CREATE2} with operator-supplied parameters
(see \S\ref{sec:admin}). The planner indexes deployed pools and tokens for
off-chain discovery (\texttt{getAllPools}, \texttt{getPoolsByToken},
\texttt{tokenIndex}). A delta-equality check on the initial deposit
(\texttt{balanceAfter $-$ balanceBefore == initialDeposits[i]}) at deploy
time rejects fee-on-transfer and rebasing tokens that fail to deliver the
exact deposit amount. Pre-existing balances at the CREATE2 address are
absorbed into reserves and gifted to the first depositor; this is intentional
and closes a deployment-griefing DoS that total-equality would have re-opened
(see \texttt{doc/security/threat-model.md} \S J.6).

\subsection{PartyInfo: View Helpers}

\texttt{PartyInfo} is a singleton view contract that accepts a target pool
address. It hosts pricing utilities (\texttt{price}, \texttt{poolPrice},
\texttt{swapAmountsForExactPrice}, quote helpers for the multi-step mint/burn
chains), the BFStore decoders (\texttt{denominators}, \texttt{fees}), and a
\texttt{working(pool)} health check. Quote helpers run the same simulation
as the on-chain entry points; their output is suitable for direct submission
to \texttt{swap}/\texttt{swapMint}/\texttt{burnSwap}.

\subsection{PartyConcierge: Router}
\label{sec:concierge}

\texttt{PartyConcierge} is the user-facing router. It accepts token addresses
rather than indices so wallets can render trades with human-readable token
names (EIP-7730 clear-signing compatibility). A single ERC-20 approval to
the Concierge covers every pool on the same chain. Funding is callback-based:
the Concierge stores call context in EIP-1153 transient storage
(\texttt{cbUser}, \texttt{cbPool}, \texttt{cbMode}, \texttt{cbEthBudget}),
validates \texttt{msg.sender == cbPool} on every callback, and refunds
residual ETH via \texttt{sweepEth}. Adds approximately 10k gas to a typical
swap.

The Concierge supports three input paths: standard approval pull,
Permit2 signature-based pull (no Concierge approval required, only the
universal Permit2 approval), and native-ETH wrap from \texttt{msg.value}.

\subsubsection{Mint Queue and Keeper Mechanism}
\label{sec:mint-queue}

Because the pool's per-window rate limit (\S\ref{sec:rate-limit}) can
force a \texttt{mint} or \texttt{swapMint} to fill only partially --- or
to revert outright when the window budget is momentarily exhausted, the
gate is transiently tripped, or the market is moving --- the Concierge
offers an \emph{optional} queue so a user can submit once and have the
remainder filled automatically as budget frees up in later blocks. The
queue is a convenience layer on top of the synchronous pool entry
points; it is opt-in (a \texttt{useQueue} flag) and requires the caller
to allow partial fills.

\paragraph{Try-first, then enqueue.} A queued request first attempts a
synchronous fill in the submitting transaction. If it fully fills,
nothing is queued. If it partially fills, or reverts with one of the
recoverable operational reverts (rate limited, volatile market,
slippage, too-small, lock-list-full), the unfilled remainder is appended
to a per-pool FIFO queue; non-recoverable reverts (deadline, killed
pool, and the like) bubble up unchanged and queue nothing.

\paragraph{No cross-block token escrow.} Consistent with the
protocol-wide rule that the system never custodies a user's tokens
across blocks, a queued request holds \emph{no} ERC-20 tokens. The
queue entry records the requester, recipient, target pool, remaining LP
to mint, a single \emph{output tolerance} $\tau_s$
(\texttt{tolerancePpm}), and a deadline; the tokens stay in the
requester's wallet under their existing approval and are pulled only at
the moment a keeper executes a fill. Crucially the entry stores a
\emph{tolerance}, not fixed per-token \texttt{maxAmountsIn} caps:
because a proportional basket mint is price-neutral, the keeper
re-derives the exact proportional draw caps from \emph{live} reserves
each block and reserves the slippage intent on the LP-output floor,
$\mathit{minLpOut} = \mathit{lpFill}\cdot(1 - \tau_s)$. A request whose
owner has since moved or de-approved the tokens (or whose live fill
would breach $\tau_s$) simply fails its execution attempt and is
requeued or, past its deadline / timeout, dropped --- no funds can be
stranded in the router.

\paragraph{Queue funding: allowance or Permit2.} A queued request may be
funded two ways, selected by a per-request flag. \emph{(i)} A standing
ERC-20 allowance to the Concierge: each keeper tranche pulls its draw
under the user's existing approval. \emph{(ii)} A standing Permit2
\texttt{AllowanceTransfer} allowance: dedicated \texttt{mint\-With\-Queue\-Permit2\-Allowance}
and \texttt{swapMint\-With\-Queue\-Permit2\-Allowance} entry points
register the allowance once at enqueue time (the user signs a single
Permit2 \texttt{PermitBatch} / \texttt{PermitSingle}), and each
subsequent keeper tranche draws against it via \texttt{transferFrom} with
no further signature. Either way the funds never leave the requester's
wallet until a tranche actually executes.

\paragraph{Fragmentation into rate-limit tranches.} A single large
queued request is filled across many blocks in \emph{tranches}, never as
one all-or-nothing fill. On each \texttt{executeMints} call the keeper
sizes the tranche to the pool's \emph{current} per-block $\gamma$-window
budget (\S\ref{sec:rate-limit}): it replicates the pool's
$\gamma$-accumulator decay off-chain to find the largest LP amount that
fits the remaining budget this block, draws exactly that, and decrements
the request's remaining LP. The keeper recomputes the proportional caps
and the $\mathit{minLpOut}$ floor from the stored tolerance $\tau_s$
against live reserves for \emph{every} tranche, so each piece
independently honors the user's slippage tolerance and a request that
spans a price move is never filled at a stale ratio. Successive tranches
drain the request as budget reopens, until it fully fills, hits its
deadline, or persistently fails slippage for
\texttt{SLIPPAGE\_TIMEOUT\_BLOCKS} (default $300$) and is canceled.

\paragraph{Permissionless keepers and the two-part fee.} Any address may
call \texttt{executeMints(pool, maxCount)} to drive queued requests from
the head, executing each as a tranche against the live pool until the
window budget is spent. Keepers are compensated by two immutable
Concierge-level fees, set at the router's construction and applied
uniformly across every pool it serves:
\begin{itemize}
  \item A small fixed \textbf{native keeper fee}
        (\texttt{NATIVE\_KEEPER\_FEE}), denominated in the chain's native
        asset and sized to cover a keeper's execution gas in the
        worst supported case (a proportional mint into a large basket)
        plus a slim premium. The submitting user attaches exactly this
        amount as \texttt{msg.value}; it is held as native escrow (the
        only value the router ever holds across blocks) and paid to the
        keeper that drives the request to a terminal state. A request
        that fully fills synchronously refunds it immediately.
  \item A \textbf{percentage keeper fee} of \textbf{$0.10\%$}
        (\texttt{KEEPER\_FEE\_PPM} $= 1000$ ppm), skimmed from each token
        the fill actually consumes and paid to the executing keeper. Its
        headroom is reserved by shrinking the tranche's LP target so the
        pool draw plus the post-draw skim always fits the requester's
        balance and allowance --- the FIFO head can never be wedged by a
        request that cannot cover the skim. This is the ongoing incentive
        that makes draining the queue profitable for keepers as soon as
        rate-limit budget reopens.
\end{itemize}

\paragraph{Cancellation and expiry.} A requester may cancel their own
queued request. If the target pool has been killed, the native escrow
is refunded to them directly; on a live pool the entry is replaced by a
tombstone and the escrow is forfeited to the next keeper that sweeps it
(an anti-griefing measure, so a flood of submit-then-cancel requests
cannot make keepers work for free). Keeper-side termination on an
expired deadline or a stale, persistently-unfillable request likewise
pays the native escrow to the terminating keeper. Queued mints carry no
keeper or queue machinery on native-asset baskets, which are excluded
from the queue because keeper execution is non-payable and cannot
auto-wrap ETH.

\section{Admin Powers and Parameter Fixity}
\label{sec:admin}

The deployment parameter tuple
\[
  (\kappa,\; \mathbf{f} = (f_0, \dots, f_{n-1}),\; \phi,\;
   \tau,\; s,\; \Gamma_{\max},\; \ell,\;
   \boldsymbol{\beta} = (\beta_0, \dots, \beta_{n-1}))
\]
is captured at \texttt{newPool} and is immutable for the lifetime of
the pool. The wedge-defense parameters
$(\tau, s, \Gamma_{\max})$ correspond to
\texttt{MINT\_DEVIATION\_PPM}, \texttt{EMA\_SHIFT\_BLOCKS}, and
\texttt{MAX\_GAMMA\_PER\_WINDOW\_PPM} respectively, and $\ell$ to the
post-mint LP-lock window \texttt{MINT\_LOCK\_BLOCKS}
(\S\ref{sec:mint-lock}); their semantics are
specified in \S\ref{sec:wedge}. The planner stores \emph{no} default
parameters of its own: the entire configuration vector above is
supplied by the pool creator at \texttt{newPool} (in the
\texttt{PoolImmutables} argument, together with $\kappa$, the fees, and
the protocol-fee settings) and captured as on-pool immutables --- a
deployed pool's parameters do not change. The
deployer enforces hard caps at construction time:

\begin{table}[h]
\centering
\caption{Hard caps enforced at \texttt{PartyPlanner.newPool}}
\label{tab:caps}
\begin{tabular}{ll}
\toprule
Parameter & Bound \\
\midrule
Per-asset swap fee $f_i$       & $< 10{,}000$ ppm ($< 1\%$) \\
Protocol fee share $\phi$      & $< 300{,}000$ ppm ($< 30\%$) \\
Asset count $n$                & $2 \le n \le 383$ \\
$\kappa$                       & $> 0$ (Q64.64) \\
Per-token base $\beta_i$       & $> 0$ \\
$\tau$ (\texttt{MINT\_DEVIATION\_PPM}) & $< 1{,}000{,}000$ ppm ($< 100\%$) \\
$s$  (\texttt{EMA\_SHIFT\_BLOCKS})     & $1 \le s < 64$ \\
$\Gamma_{\max}$ (\texttt{MAX\_GAMMA\_PER\_WINDOW\_PPM}) & $> 0$ \\
$\ell$ (\texttt{MINT\_LOCK\_BLOCKS}) & $\le 50{,}400$ (one week of L1 blocks) \\
\bottomrule
\end{tabular}
\end{table}

\subsection{Mutable State on a Live Pool}

The only mutable state on a deployed pool, post-\texttt{newPool}, is:

\begin{itemize}
  \item \texttt{protocolFeeAddress} --- recipient of accrued protocol fees,
        settable by the pool owner via \texttt{setProtocolFeeAddress}. Zero
        address is rejected when $\phi > 0$. The recipient is read at
        \texttt{collectProtocolFees} call time (see \S\ref{sec:fees}).
  \item Two-step ownership handoff via the \texttt{Ownable2Step} surface.
  \item The irreversible \texttt{kill()} flag.
\end{itemize}

\subsection{Emergency Kill Switch}

The pool owner may call \texttt{kill()}, after which the following
state-mutating entry points revert: \texttt{swap}, \texttt{mint},
\texttt{swapMint}, \texttt{burnSwap}, and \texttt{initialMint}. The
action is one-way; there is no \texttt{unkill}.
Both \texttt{burn} (proportional LP redemption) and
\texttt{collectProtocolFees} (sweep of already-accrued protocol fees to
the configured \texttt{protocolFeeAddress}) remain operational on a
killed pool by design: the proportional-redemption guarantee preserves
LP exit under any kill scenario, and continued availability of
\texttt{collectProtocolFees} ensures already-accrued protocol-fee
balances are not stranded. \texttt{kill()} is intended for incident
response if a critical vulnerability is discovered post-deployment; it is
disjoint from the immutable parameter set.

\subsection{What the Operator Cannot Do}

The operator (planner owner and pool owner) explicitly \emph{cannot}:
mint LP out of band, burn user LP, change $\kappa$ or any fee parameter on
a live pool, pause user-side burns, blacklist or seize user balances, upgrade
the pool implementation, renounce ownership, or move funds out of the pool
directly. The full inventory of admin powers (CAN and CANNOT) is documented
in \texttt{doc/security/admin-powers.md}.

\section{Deployment Guidelines}
\label{sec:deploy}

\subsection{Operator Workflow}

The recommended end-to-end deployment workflow is:

\begin{enumerate}
  \item \textbf{Token vetting.} Run \texttt{bin/validate-token <addr>}
        against every candidate token (see
        \texttt{doc/\allowbreak security/\allowbreak trusted-deployer-policy.md}
        for the validator policy). Resolve every \texttt{WARN} off-chain;
        \texttt{FAIL} blocks listing. Cross-check with
        \texttt{slither-check-erc <addr> ERC20 --erc-conformance}. Record
        verifications in a listing record.
  \item \textbf{Token-set composition.} Choose the basket of assets so the
        pool sits in a low-LVR regime
        (\S\ref{sec:token-selection})~--- prefer correlated assets, group
        natural baskets, and avoid mixing high-volatility singletons with
        otherwise correlated sets.
  \item \textbf{Parameter selection.} Choose $\kappa$ (\S\ref{sec:kappa}),
        per-asset fees, the $\sigma_{\mathrm{swap}}$ gate triple
        $(\tau, s, \Gamma_{\max})$ (per the per-family table in
        \texttt{doc/rate-limited-mints.md}), and bases
        (\S\ref{sec:bases}).
  \item \textbf{Initial liquidity.} Fund the deployer with
        \texttt{initialDeposits[i]} for each token; approve the planner.
  \item \textbf{Deploy.} Call \texttt{PartyPlanner.newPool(...)}. The
        planner deploys the pool via \texttt{CREATE2}, pulls initial
        liquidity from the operator-specified \texttt{payer}, and invokes
        \texttt{initialMint} on the new pool, minting the operator-chosen
        initial LP supply $L^{(0)}$ (defaulting to $10^{18}$ if zero is
        passed) to \texttt{receiver}.
  \item \textbf{Verify.} Confirm the pool address, \texttt{LMSR()} state, and
        \texttt{balances()} match expectations. Subscribe to the
        \texttt{PartyStarted} event for the indexer.
\end{enumerate}

\subsection{Token-Selection Criteria: Correlated-Asset LVR Reduction}
\label{sec:token-selection}

The operator's choice of which tokens to bundle into a single pool is a
first-class design lever, separate from the per-asset fee and $\kappa$
calibration. The load-bearing criterion is correlation structure across
the basket, because LMSR's multi-asset geometry exhibits a
\emph{cross-correlation discount} on the loss-versus-rebalancing (LVR)
rate that is absent from separable bilateral CFMMs.

\paragraph{The discount.}
Olson \& Claude Opus (2026, in-tree
\texttt{doc/lvr-multiasset-lmsr.pdf}) derive the instantaneous LVR rate
for the Hanson multi-asset LMSR in closed form. Writing $P_i$ for
external prices of $N$ risky assets, $S = \sum_i P_i$, $b$ for the LMSR
liquidity parameter, and $\Sigma$ for the return-covariance matrix, the
rate decomposes as
\[
  \ell(P)
    \;=\;
    \underbrace{\tfrac{b}{2}\sum_i \sigma_i^2 P_i}_{\text{$N$ independent two-asset LMSR rates}}
    \;-\;
    \underbrace{\tfrac{b}{2(1 + S)} P^{\!\top}\!\Sigma P}_{\text{cross-correlation discount}}.
\]
The discount is non-negative for any positive-semidefinite $\Sigma$ and
grows with the variance of the value-weighted basket return: the more
correlated the basket, the more cross-correlation cancels relative-mode
arbitrage. For perfectly correlated equal-volatility assets the
relative-mode contribution vanishes entirely and an $(N{+}1)$-asset
balanced pool collapses to the LVR of a single two-asset LMSR with
matched $b$ and matched risky notional --- the multi-asset construction
buys $N(N{-}1)/2$ pairwise quotes at the LVR cost of one. A pool of
perfectly correlated assets with no num\'eraire has zero LVR.

\paragraph{Operator implication.}
Token selection should treat correlation as a primary criterion, not a
secondary one. Concretely:
\begin{itemize}
  \item \textbf{Prefer correlated baskets.} A liquid-staking-ETH bundle
        (stETH, rETH, ETH), a stablecoin bundle (USDC, USDT, DAI), a
        wrapped-BTC bundle (wBTC, cbBTC, tBTC), or any
        sector-correlated long-tail set will retain more value against
        the same arbitrage pressure than the same tokens spread across
        disjoint pairwise pools.
  \item \textbf{Avoid mixing.} A single highly-volatile uncorrelated
        asset inserted into an otherwise correlated basket re-introduces
        the relative-mode LVR contribution between that asset and every
        other slot; it is usually better to deploy the outlier in its
        own pool (or pair) than to dilute the correlation benefit of the
        surrounding basket.
  \item \textbf{Group by num\'eraire vs.\ risky.} A no-num\'eraire pool
        of correlated assets is the extreme of the discount; if a
        num\'eraire is included, the common-mode num\'eraire term
        dominates and the LVR floor is set by the basket-vs-num\'eraire
        variance.
\end{itemize}
The LVR paper's closed-form decomposition gives the operator a
quantitative tool to compare candidate baskets before deploying: compute
$P^{\!\top}\!\Sigma P / (\sum_i \sigma_i^2 P_i)$ on a candidate
composition, larger is better. The same calculation is the input to the
$\kappa$ and $C$ calibration steps in
\S\ref{sec:kappa}--\S\ref{sec:fee-calibration}.

\subsection{Choosing $\kappa$}
\label{sec:kappa}

$\kappa$ sets the liquidity parameter $b = \kappa\,\Sm$ and so controls
the trade-off between depth and drain-resistance. A higher $\kappa$
deepens the pool --- flatter price curve, tighter quotes, less slippage
per trade --- but lowers the external single-asset price move the pool
absorbs before that slot drains; a lower $\kappa$ steepens the curve so
the pool re-prices through larger moves without draining. The maximum
external single-asset price move before that slot drains is
\begin{equation}
  R_{\mathrm{p,max}} \;=\; \exp\!\left(\frac{1}{\kappa\,(n-1)}\right).
\end{equation}
For multi-asset pools the same $\kappa\,\Sm$ liquidity budget is shared
across $n$ slots, so single-slot divergence drains faster as $n$ grows.

\paragraph{Calibration is simulation-driven, not table-driven.}
$\kappa$ is not read off a generic by-asset-class table. It is tuned
per pool with an off-chain analyzer (not part of the public repository)
that simulates arbitrage actors against historical solver order flow
(CoW) on the pool's \emph{fixed} token composition and fee schedule,
sweeping $\kappa$ for the value that maximizes LP yield without draining
all but one slot. Because the fee schedule is an input to that
simulation, \textbf{fees are chosen first} (\S\ref{sec:fee-calibration})
and $\kappa$ is tuned against them. The two production pools land far
apart: the volatile OG basket uses $\kappa = 0.2$ (a steep curve that
survives large blue-chip and long-tail moves), while the correlated
Peg Party stablecoin basket uses $\kappa = 5$ (a deep, tightly-quoting
curve that drains a slot only on a genuine depeg --- the desired
``toxic flow stays out'' behavior for assets assumed to hold their
peg). See \texttt{analyzer/og-pool.md} and
\texttt{analyzer/stable-pool.md}.

\subsection{Bases $\beta_i$ and Market-Calibrated Initialization}
\label{sec:bases}

Each per-token base $\beta_i$ converts raw token units to internal Q64.64
units: $q_i = \mathrm{balance}_i / \beta_i$. The bases are not chosen
freely --- they are fixed by the \texttt{initialDeposits} array passed to
\texttt{PartyPlanner.newPool}, in which $\beta_i$ equals the per-token
deposit amount that the deployer commits at \texttt{initialMint}. Because
the kernel's internal coordinate is $q_i = \mathrm{balance}_i / \beta_i$,
this means \emph{every slot initializes at $q_i = 1.0$} regardless of the
underlying token's raw-unit magnitude or USD price.

\paragraph{Market-calibrated initial deposits.}
The deployment script that prepares the \texttt{initialDeposits} array
is external to the pool repository and operates as follows:

\begin{enumerate}
  \item Query realtime market prices for every candidate asset at
        deployment time from an off-chain reference (centralized exchange
        feed, oracle aggregator, etc.).
  \item Compute a per-token initial deposit amount such that the
        \emph{USD value} of each token's initial deposit is equal across
        all slots in the new pool.
  \item Pass that array to \texttt{newPool} as \texttt{initialDeposits};
        the planner forwards it as the pool's immutable $\boldsymbol{\beta}$
        vector and pulls the matching amounts at \texttt{initialMint}.
\end{enumerate}

The result is a pool whose initial internal state has
$q_i = 1.0$ for every $i$, with the internal $1{:}1$ marginal price
between any two slots matching the slots' realtime external market price
ratio at deployment. The pool is launched in price-equilibrium with the
external market.

\paragraph{Mis-calibration is an immediate-arbitrage failure mode.}
If the deployer's initial-deposit ratios do \emph{not} match the realtime
external prices at deployment, the pool launches in disequilibrium: the
internal $1{:}1$ marginal price differs from the external market price,
and arbitrageurs can extract the resulting spread immediately at the
initial depositor's expense. The pool itself has no internal way to
detect this (it has no oracle); enforcement is entirely a property of
the deployment script. The operator's discipline is to run the
calibration script as close to the deployment transaction as practical
and to abort if intermediate price drift exceeds a small tolerance.

\paragraph{Why this works.}
Equal-USD-value deposits give every $q_i$ identical magnitude in internal
units, which:

\begin{itemize}
  \item Maximizes precision headroom across natural drift.
  \item Aligns the spot price near $1{:}1$ in internal units (true at the
        instant of launch, drifting thereafter with external prices).
  \item Sits the pool well inside the kernel's $\mathrm{EXP\_LIMIT}$
        envelope.
\end{itemize}

Bases are immutable; revising them requires deploying a new pool.

\subsection{Fee Calibration}
\label{sec:fee-calibration}

Per-asset fees use a \textbf{$C \cdot \sigma$ method}: within a given pool,
the fee rate for asset $i$ is set proportional to a volatility estimate
$\sigma_i$ for that asset, scaled by a per-pool constant $C$. USDC
anchors the schedule at $1.0$ bps; every other listed asset's fee is set
so that $f_i / \sigma_i$ matches USDC's $f_{\mathrm{USDC}} / \sigma_{\mathrm{USDC}}$,
then scaled by $C$. The scaling constant $C$ is \emph{not} protocol-wide
--- it is chosen per pool from competitive analysis of the venues that
quote the same or comparable assets.

\paragraph{Bridge surcharge.}
Wrapped-bridged assets (e.g.\ \texttt{wSOL}) carry an additional
heuristically-applied fee component on top of the $C\sigma$ baseline. The
surcharge reflects the bridge-trust risk premium: an LP holding a wrapped
asset bears the bridge's solvency and operational risk in addition to the
underlying asset's market risk, and the fee schedule compensates them for
that.

\paragraph{Step 1: calibrate $C$ from per-pool competitive analysis.}
$C$ is set as a first step, before any $\kappa$ tuning, by reference to
the fees charged by the largest-TVL competing venues on the pool's asset
composition. $C$ is raised to the largest value that still ``just
undercuts'' the prevailing market rate on the relevant pairs --- high
enough to extract meaningful fee revenue from organic flow, low enough
that fee-sensitive takers prefer Liquidity Party to the competing venue.
Once chosen, $C$ is fixed for the rest of the calibration of \emph{this
pool}; a different pool composition may select a different $C$.

\paragraph{Reading the competing venue's true market rate.}
Several mainstream AMMs (Uniswap v3 is the canonical case) expose only a
coarse set of fee tiers --- typically $\{1, 5, 30, 100\}$ bps --- so the
tier a competing pool sits on is not always a clean signal of the
prevailing market rate. The operator's job is to recover the
\emph{economic} market fee from the deployed tier by judging the tier's
TVL share against the competing tiers on the same pair, and by backing
out any structural cost that the tier price covers but that does not
apply to a marginal Liquidity Party depositor. The canonical example is
WBTC/USDT, which clears at the 30~bps tier on Uniswap v3 despite the
underlying asset volatility implying a much smaller economically-fair
rate; the operator interprets the residual as the LP's structural cost
of the WBTC redemption window (roughly $\sim 20$~bps of one-way
redemption friction that WBTC LPs must price in to avoid losing on
forced-redemption events). That structural cost is absent for a marginal
Liquidity Party LP, so the operator subtracts it discretionarily to
arrive at the volatility-matched market rate used in the $C\cdot\sigma$
fit for the new pool. Similar discretionary adjustments apply to any
competing venue whose deployed fee covers redemption, custody, or other
structural LP costs that do not transfer to a Liquidity Party deposit.

\paragraph{Step 2: optimize $\kappa$ at fixed $C$.}
With $C$ fixed, $\kappa$ is tuned against a historical simulation over
at least one year of price and flow data. The simulation operates two flow streams against a candidate
pool:

\begin{enumerate}
  \item A \emph{Binance-connected arbitrageur} able to take infinite
        liquidity at the Binance inside quote plus 2 bps. This is an
        upper bound on real arbitrage activity --- the inside is thin in
        practice, so a real arb cannot take infinite size at the inside
        plus 2 bps. The simulation therefore over-estimates arbitrage
        pressure on the pool, biasing $\kappa$ toward conservative
        (LP-favoring) values.
  \item \emph{Organic flow} derived from on-chain trade data, specifically
        CoW Swap auction participant flow over the same historical window.
        This is real, fee-sensitive user demand.
\end{enumerate}

The objective is the LP's terminal net return: fees collected minus
LVR-style adverse-selection cost minus realized IL on slot drift, measured
at the simulated horizon. $\kappa$ is chosen to maximize that net at the
fixed $C$.

The procedure --- $C$ from per-pool competitive analysis, then $\kappa$
from LP-net-return optimization at the chosen $C$ --- decouples the two
parameters: $C$ reflects the pool's competitive position and $\kappa$
reflects historical LP economics under that competitive position. Both
are re-run for each new pool whose asset composition or competitive
landscape materially differs from a previously calibrated pool; reuse is
allowed only when the new pool genuinely sits in the same competitive
envelope.

\subsection{Native-Asset Wrapping}

If the pool will accept native ETH (or chain-native equivalent), one of the
pool's tokens must be the canonical wrapper (\texttt{WETH9} on mainnet, the
equivalent on each L2). The wrapper address is passed via the planner
constructor and is the only address authorized to push raw ETH into the pool
through \texttt{receive()}.

\section{Security Review Process}
\label{sec:security}

\subsection{Layered Defense}

The protocol applies a defense-in-depth posture
(\texttt{doc/\allowbreak security/\allowbreak threat-model.md} \S1):

\begin{itemize}
  \item \textbf{Permissioned deploy} as the primary defense against
        malicious tokens. \texttt{PartyPlanner.\allowbreak newPool} is
        \texttt{onlyOwner}, indefinitely. The operator is the trust
        anchor for token vetting.
  \item \textbf{Runtime guards.} \texttt{nonReentrant} on every state-mutating
        entry point, \texttt{killable} on every entry point except
        \texttt{burn} and \texttt{collectProtocolFees}, payer-gates on the
        \texttt{APPROVAL} and
        \texttt{PREFUNDING} funding paths, deploy-time delta-equality on the
        initial-deposit transfer, EIP-712 witness binding on every Permit2
        path.
  \item \textbf{Off-chain validation.} \texttt{bin/validate-token} +
        \texttt{slither-check-erc} are mandatory pre-listing. The validator
        emits \texttt{PASS}/\texttt{WARN}/\texttt{FAIL}; \texttt{WARN}
        requires off-chain verification and a recorded listing-record entry,
        \texttt{FAIL} blocks the listing entirely.
  \item \textbf{Monitoring + kill switch.} Off-chain drift monitoring with
        documented yellow/red thresholds; \texttt{kill()} as the last-resort
        lever.
\end{itemize}

There is no on-chain price oracle dependency; pricing is intrinsic to the
LMSR kernel.

\subsection{Internal Review Artifacts}

The repository carries the following internal review artifacts that gate
launch (see \texttt{doc/security/}):

\begin{itemize}
  \item \texttt{threat-model.md} --- threat model, actors, fund-bearing
        assets, attack vectors, invariants. v1 pre-deploy.
  \item \texttt{checklist.md} --- the security review checklist (14
        sections, 84 rows). \textbf{Every row is closed.} Engineering and
        code-correctness rows are recorded as \textsc{mitigated} (closed
        in source) or \textsc{documented} (closed by docs); the external
        audit is complete (\S\ref{sec:external-audit}), leaving one
        operational row --- the migration to a multisig admin key ---
        recorded as \textsc{deferred with trigger} at a named TVL
        threshold, with compensating mitigations specified in
        \texttt{threat-model.md} \S 11. No row is open or
        unresolved.
  \item \texttt{asset-authority-matrix.md} --- (actor $\times$ asset $\times$
        function) $\to$ gate $@$ file:line, the per-cell authorization matrix.
  \item \texttt{admin-powers.md} --- CAN/CANNOT inventory of every
        \texttt{onlyOwner} surface, with trapped-funds analysis.
  \item \texttt{trusted-deployer-policy.md} --- operator obligations, the
        token-class denylist (fee-on-transfer, rebasing, ERC-777-style hook
        tokens, multi-address proxy, mutable governance), the listing
        checklist.
  \item \texttt{economics-audit.md} --- the economics-side review of the
        fee structure, single-asset round trips, and LP-favor invariants.
  \item \texttt{storage-layout-audit.md} --- audit of the C3-linearized
        storage layout used under \texttt{DELEGATECALL}.
  \item \texttt{unchecked-blocks-audit.md} --- justification for every
        \texttt{unchecked} block in \texttt{src/}.
  \item \texttt{differential-review.md} --- differential review of the
        block-anchored $b$ kernel and the wedge defense.
  \item \texttt{tooling-runbook.md} --- Slither, Mythril, and forge-test
        invocation reference.
  \item \texttt{user-allowance-guidance.md} --- guidance for integrators on
        the Concierge approval surface and Permit2 usage.
\end{itemize}

\subsection{Tooling}

\begin{itemize}
  \item \textbf{Foundry / forge-test.} The in-tree test suite carries
        the cost-preservation, kernel-no-overshoot, gate-tripping,
        rate-limit-boundary, min-clamped-burn, JIT-closure,
        wedge-attack, and (actor $\times$ asset) gating tests that
        close the checklist.
        Coverage is collected with
        \texttt{forge coverage --ir-minimum --no-match-contract "ContractSize|GasTest"}.
  \item \textbf{Slither.} Static analysis with the project's
        \texttt{slither.config.json}. False-positive suppressions are
        committed alongside the code they suppress and are reviewed in the
        checklist.
  \item \textbf{Token validator.} \texttt{bin/validate-token} probes a
        candidate token across the property axes that the protocol cannot
        defend against at runtime (fee-on-transfer, rebasing, hook
        callbacks, multi-address proxy, mutable governance). See
        \texttt{doc/\allowbreak security/\allowbreak trusted-deployer-policy.md}.
\end{itemize}

\subsection{External Audit}
\label{sec:external-audit}

The codebase was audited in its entirety over 15 full working days by
\href{https://github.com/jayjoshix}{Jay Joshi}. All findings were
mitigated.

\subsection{Bug Bounty}

Liquidity Party is launching as a pre-revenue protocol and is not in a
position to commit to a formal bug bounty program prior to launch. Per
\texttt{SECURITY.md}, severity-graded discretionary thanks are extended for
responsible reports during the interim period. The intent is to launch a
formal bounty (Immunefi Boost or equivalent) once launch metrics justify the
spend. This section will be updated once a bounty program has been
established.

\subsection{Vulnerability Disclosure}

Security reports are routed privately by email to
\texttt{tim@\allowbreak dexorder.com} or via direct message to
\href{https://x.com/LiquidityParty}{\texttt{@LiquidityParty}} on X.
Public GitHub issues for security vulnerabilities are explicitly
discouraged. The disclosure SLA is acknowledgement within 48 hours and
a coordinated disclosure timeline scaled to severity. The reporter is
credited at their option once the issue is mitigated. The full policy
lives in \texttt{SECURITY.md}.

\subsection{Out of Scope}

Per \texttt{SECURITY.md}, the following are out of scope:

\begin{itemize}
  \item Test fixtures, mocks, and helpers under \texttt{test/}.
  \item Issues that require a compromised admin key as a precondition. Admin
        keys can \texttt{kill()} a pool and redirect protocol fees, but
        cannot reach LP reserves; this is documented as accepted risk in
        \texttt{doc/security/threat-model.md} \S11.N7 and
        \texttt{doc/security/admin-powers.md}.
\end{itemize}

\section{Risk Management}
\label{sec:risk}

\subsection{Token-Slot Drainage Is Expected LP Behavior}
\label{sec:drainage}

\textbf{Token slots in a Liquidity Party pool will drain. This is normal,
expected, and intentional.} (We use ``token slot'' for a pool's holding
of one of its constituent tokens --- the position the pool maintains in,
for example, the WBTC token within a WBTC/USDC/ETH pool.) The protocol's
deployment strategy is to calibrate $(C, \kappa)$ for maximum LP net
return after impermanent loss (\S\ref{sec:fee-calibration}). Optimizing
for net LP return implies that token slots whose underlying asset
experiences a large persistent price move \emph{will} reach the drained
state $q_k \approx 0$ before the pool internalizes the full move. This is
the same trade-off every concentrated-liquidity AMM makes.

\paragraph{Analogy: Uniswap v3 range positions.}
A Uniswap v3 LP who posts liquidity over a finite price range collects
fees only while the spot price is inside the range. If the price moves
through the upper end of the range, the LP ends up holding 100\% of the
``cheaper'' asset --- their position has converted entirely to the
appreciated asset's pair-leg, exactly as if they had executed a market
sell at every intermediate price along the range. Range LPs accept this
as the fundamental staking trade-off of v3: deeper liquidity inside the
range in exchange for full conversion at the boundary. Drainage in
Liquidity Party is the multi-asset analogue. A token slot is implicitly
``range-staked'' against the rest of the basket: while the relative price
stays within the pool's depth envelope, the slot earns fees on the LMSR
quote; if the external market moves the slot's price far enough that the
pool can fund cheaper-than-external swaps until the slot empties,
arbitrageurs will drain it. The maximum single-token-slot external move
the pool can absorb before drain is
\begin{equation}
  R_{\mathrm{p,max}} \;=\; \exp\!\left(\frac{1}{\kappa\,(n - 1)}\right)
\end{equation}
(\S\ref{sec:kappa}); past that, the token slot empties and the basket
settles in the LP's hands without that asset.

\paragraph{What the LP retains after a drain.}
A drained token slot does not impair the LP's redemption rights.
Proportional \texttt{burn} continues to work and returns the LP their
share of every token slot that still has reserves. Single-asset
\texttt{burnSwap} into the drained token returns nothing for that token
but treats the drained slot as zero-contribution rather than reverting;
\texttt{burnSwap} into a non-drained token still works. The LP's economic
position after a token-slot drain is equivalent to holding the surviving
basket --- exactly the position a v3 LP holds after a one-sided range
exit.

\paragraph{Industry framing.}
Drainage in this sense is the industry-standard cost of providing
concentrated liquidity. Liquidity Party does not promise to retain every
token under every market scenario; it promises (a) pro-rata redemption
of whatever the pool currently holds, (b) LP-favoring rounding and
approximation at every step, and (c) a fee schedule calibrated against
historical data with the goal of anticipated LP net return inclusive of
expected drainage events. The historical calibration is a best-effort
estimate; past flow and price dynamics are not a guarantee of future LP
returns, and the realized net is subject to the market conditions that
unfold after deployment. LPs participating in pools that hold
uncorrelated or highly volatile assets should expect drain events on
individual token slots and should size their participation accordingly.

\subsection{LP Loss Bound}

Under constant $b$, classical LMSR admits a worst-case loss bound of
$b\,\ln n$ in the payoff num\'eraire. With $\bofq = \kappa\,\Sm$, the
per-state bound scales with the current size metric: instantaneous worst-case
loss $\bofq\,\ln n = \kappa\,\Sm\,\ln n$, or $\kappa\,\ln n$ per unit of
$\Sm$. Because $b$ varies with state, global worst-case loss along an
arbitrary path depends on how $\Sm$ evolves; the proportionality clarifies
how risk scales with deployed liquidity.

\subsection{Liquidity-Manipulation Resistance}

Every operation that changes $\Sm$ is priced through the same LMSR potential
that prices ordinary swaps:

\begin{itemize}
  \item Proportional mint and burn are exact scalings:
        $\mathbf{q} \mapsto (1+\alpha)\mathbf{q}$ implies
        $b \mapsto (1+\alpha)b$ and every pairwise marginal ratio is
        invariant. A proportional mint followed by its exact inverse returns
        the pool to the same state with no net transfer.
  \item Single-asset mint and burn (\texttt{swapMint}, \texttt{burnSwap})
        decompose into a proportional rescaling composed with a bundle of
        LMSR swaps priced by the same convex potential. The cost of changing
        $\Sm$ is the sum of kernel-priced components; there is no off-market
        ``free'' lever to change $b$.
\end{itemize}

Because $C(\mathbf{q})$ is convex, the net kernel cost of any closed
cycle of allowed operations sums to zero in the fee-free idealization;
any practical implementation cost (fees, LP-favor rounding) makes the
cycle strictly loss-making for the attacker. The block-anchored $b$
(\S\ref{sec:block-b}) ensures that every leg of every intra-block
sequence of \texttt{swap}, \texttt{swapMint}, and \texttt{burnSwap}
prices against the same convex potential; combined with the
wedge defense (\S\ref{sec:wedge}) on the mint side,
adversarial closed cycles run on already-imbalanced pools cannot
extract a single wei (\S\ref{sec:lp-safety}, ``Empirical
zero-extraction'').

\subsection{Numerical Envelope}

The kernel's cost-preservation residual sits below $10^{-9}$ relative across
the production-relevant range of pool imbalances ($R_q \le 10^7$) and degrades
gracefully past that. The ABDK exp/ln round-trip noise floor is
$\sim 2 \times 10^{-16}$ relative; production tolerances are bound at
$10^{-9}$, leaving seven orders of magnitude of headroom (see
\texttt{security/drainage-guidelines.md} \S3, \S4).

\subsection{Operational Mitigations}

\begin{itemize}
  \item User \texttt{minAmountOut} / \texttt{maxAmountIn} bounds on every
        state-mutating entry point.
  \item Capacity caps on outputs ($y \le q_j$), enforced inside the kernel.
  \item Minimum-balance constraints on \texttt{initialMint} (every initial
        deposit must be strictly positive; \texttt{$\beta_i > 0$} enforced at
        deploy time).
  \item Off-chain drift monitoring with yellow/red thresholds on the
        token-slot drain ratio $\min q_i / (\Sm/n)$ and the q-ratio
        $\max q / \min q$.
  \item \texttt{kill()} as the last-resort lever, preserving proportional
        burn for LP exit.
\end{itemize}

\subsection{Out-of-Scope Token Classes}

The protocol does not support, and the trusted-deployer policy MUST NOT
list:

\begin{itemize}
  \item Fee-on-transfer tokens.
  \item Rebasing tokens.
  \item ERC-777-style hook tokens, ERC-677 \texttt{transferAndCall} tokens,
        any token that callbacks during \texttt{transfer} or
        \texttt{transferFrom}.
  \item Multi-address proxy tokens.
  \item Tokens whose owner can later add fees, pause transfers, or
        blacklist addresses (mutable post-list governance).
\end{itemize}

A token that the validator clears today but later mutates into one of the
above classes is a class-2 hazard; the response is \texttt{kill()}.

\section{Deployment Surface}
\label{sec:deployments}

\paragraph{Initial deployment.}
The initial production deployment targets Ethereum mainnet. The first
planned expansion, if launch metrics support it, is Base. The protocol is
EVM-only --- non-EVM chains are out of scope for the foreseeable future.
Per-chain contract addresses (\texttt{PartyPlanner}, \texttt{PartyInfo},
\texttt{PartyConcierge}) are published in
\texttt{deployment/liqp-deployments.json} in the repository. That file
is the canonical record of deployed contract information and is updated
at deploy time.

\paragraph{Operator key custody.}
The admin (planner owner / pool owner) key is held on a hardware wallet.
No further detail of the operator's internal key-management procedure is
disclosed; the LP's security guarantee does not depend on it
(\S\ref{sec:admin-risk}). The protocol is designed to be \emph{resilient
against an admin-key leak}: a leaked admin key cannot reach LP reserves
under any operation in the contract surface. The worst case under a leaked
key is denial-of-service (\texttt{kill()} called by the attacker) plus
redirection of not-yet-collected protocol fees; LP \texttt{burn} remains
operational and LPs retain their pro-rata share of reserves. The core team
retains independent ability to invoke \texttt{kill()} on the planner-owned
pools as a final defensive measure if the admin key is compromised.

\paragraph{Indexer.}
Liquidity Party does not host a public subgraph at this time. Token and
pool metadata is published as a static \texttt{deployment/metadata.json}
artifact, which is the canonical source for that metadata;
\texttt{deployment/liqp-deployments.json} holds only the deployed
contract information.

Integrators that need event-level data should consume the canonical event
set directly from the chain:
\texttt{PartyStarted}, \texttt{Mint}, \texttt{Burn}, \texttt{Swap},
\texttt{SwapMint}, \texttt{BurnSwap},
\texttt{ProtocolFeesCollected}, \texttt{Killed}.

\section{Disclaimers}
\label{sec:disclaim}

\paragraph{Issuer.}
Liquidity Party is operated by Dexorder Trading Services, Ltd., a company
incorporated in the British Virgin Islands under registration
number 2160219.

\paragraph{Access restrictions.}
The Liquidity Party front-end blocks access from addresses appearing on the
U.S. Office of Foreign Assets Control (OFAC) Specially Designated Nationals
(SDN) and consolidated sanctions lists. The on-chain smart contracts are
permissionless to interact with from any address; the access restriction is
enforced at the front-end and any sanctioned address that interacts
directly with the contracts does so outside the front-end's terms of
service.

\paragraph{No investment advice.}
This whitepaper is for informational purposes only. It does not constitute
investment advice, a solicitation of any security, or a representation as
to the suitability of the protocol for any user, jurisdiction, or
purpose. Past performance of the simulation environment described in
\S\ref{sec:fee-calibration} is not indicative of future LP returns.

\paragraph{No representation as to tax treatment.}
The tax treatment of LP participation, swap activity, and protocol-fee
distribution is jurisdiction-specific and depends on facts and
circumstances outside the protocol's control. Users are responsible
for their own tax compliance.

\paragraph{Forward-looking statements.}
Statements in this whitepaper regarding future deployments (Base
expansion, formal bug bounty, metadata index), future audit report
publication, and future protocol behavior are forward-looking. They
reflect current intent and are not commitments.

\paragraph{Risk acknowledgement.}
Liquidity provision in Liquidity Party pools is subject to impermanent
loss, token-slot drainage loss (\S\ref{sec:drainage}), smart-contract risk,
oracle-independent depegging of constituent assets, bridge risk on
wrapped-bridged assets, and the residual risks documented in
\S\ref{sec:risk}. LPs should understand these risks before depositing.

\section*{References}
\addcontentsline{toc}{section}{References}

\begin{itemize}
  \item Hanson, R. (2002). \emph{Logarithmic Market Scoring Rules for Modular
        Combinatorial Information Aggregation.}
        \href{https://mason.gmu.edu/~rhanson/mktscore.pdf}{mason.gmu.edu/\textasciitilde{}rhanson/mktscore.pdf}
  \item Othman, A., Pennock, D., Reeves, D., Sandholm, T. (2013).
        \emph{A Practical Liquidity-Sensitive Automated Market Maker.}
        \href{https://www.cs.cmu.edu/~sandholm/liquidity-sensitive\%20market\%20maker.EC10.pdf}{cs.cmu.edu/\textasciitilde{}sandholm/liquidity-sensitive market maker.EC10.pdf}
  \item Olson, T., Claude Opus 4.7 (2026). \emph{Loss-Versus-Rebalancing
        Reduced in Cross-Correlated Multi-Asset LMSR Pools.}
        \texttt{doc/lvr-multiasset-lmsr.pdf} (this repository).
\end{itemize}

\end{document}
